\element{quizz}{
  \begin{question}{Q001}
    Big Data เดิมถูกอธิบายด้วย 3 องค์ประกอบหลัก คืออะไร
    \begin{choices}
        \correctchoice{Volume, Velocity, Variety}
        \wrongchoice{Value, Validity, Veracity}
        \wrongchoice{Visualization, Verification, Variation}
        \wrongchoice{Volume, Value, Veracity}
    \end{choices}
  \end{question}
}

\element{quizz}{
  \begin{question}{Q002}
    ข้อใดไม่ใช่ความท้าทายของ Big Data ตามนิยาม
    \begin{choices}
        \wrongchoice{Capturing data}
        \wrongchoice{Data storage}
        \wrongchoice{Data analysis}
        \correctchoice{Data cooking}
    \end{choices}
  \end{question}
}

% \element{quizz}{
%   \begin{question}{Q003}
%     ตัวเลือกใด ``จับคู่ผิด
%     \begin{choices}
%         \wrongchoice{Library search = Old Fashion}
%         \wrongchoice{Trade publication extraction = Old Fashion}
%         \correctchoice{Geolocation stream = Old Fashion}
%         \wrongchoice{Social networks API = Today}
%     \end{choices}
%   \end{question}
% }

\element{quizz}{
  \begin{question}{Q004}
    ข้อใดตรงกับ Structured Data ที่สุด
    \begin{choices}
        \wrongchoice{ข้อมูลที่ไม่มีรูปแบบตายตัว}
        \correctchoice{ข้อมูลอยู่ในตาราง รูปแบบ แถว,คอลัมน์ชัดเจน}
        \wrongchoice{ข้อมูลจากไฟล์เสียง}
        \wrongchoice{ข้อมูลรูปภาพ}
    \end{choices}
  \end{question}
}

\element{quizz}{
  \begin{question}{Q005}
    ข้อใดเป็น Unstructured Data
    \begin{choices}
        \wrongchoice{ตาราง SQL}
        \correctchoice{ข้อความอิสระ, วิดีโอ, ไฟล์เสียง}
        \wrongchoice{ไฟล์ Excel}
        \wrongchoice{ข้อมูลอยู่ในตาราง รูปแบบ แถว,คอลัมน์ชัดเจน}
    \end{choices}
  \end{question}
}

\element{quizz}{
  \begin{question}{Q006}
    ข้อใด \textbf{ไม่ใช่} ตัวอย่างของ Unstructured Data
    \begin{choices}
        \wrongchoice{รูปภาพ}
        \wrongchoice{ไฟล์ Word}
        \correctchoice{ฐานข้อมูล SQL}
        \wrongchoice{คลิปเสียง}
    \end{choices}
  \end{question}
}

\element{quizz}{
  \begin{question}{Q007}
    Metadata คืออะไร
    \begin{choices}
        \wrongchoice{ข้อมูลที่ไม่สำคัญ}
        \correctchoice{ข้อมูลที่ใช้อธิบายข้อมูล}
        \wrongchoice{ข้อมูลข้อความ}
        \wrongchoice{ข้อมูลที่อยู่ในรูปแบบ JSON}
    \end{choices}
  \end{question}
}

\element{quizz}{
  \begin{question}{Q008}
    ข้อใดกล่าวถูกต้อง
    \begin{choices}
        \correctchoice{Semi-structured มีแท็กแต่ไม่เป็นตารางตายตัว}
        \wrongchoice{Structured เป็นข้อมูลที่เป็นรูปภาพ}
        \wrongchoice{Semi-structured ไม่มี metadata}
        \wrongchoice{ไม่มีความแตกต่าง}
    \end{choices}
  \end{question}
}

\element{quizz}{
  \begin{question}{Q009}
    ข้อใดอธิบาย Volume ได้ถูกต้องที่สุด
    \begin{choices}
        \correctchoice{ปริมาณข้อมูลที่ถูกสร้างและจัดเก็บ}
        \wrongchoice{ความหลากหลายของชนิดข้อมูล}
        \wrongchoice{ความเร็วในการสร้างข้อมูล}
        \wrongchoice{ความถูกต้องแม่นยำของข้อมูล}
    \end{choices}
  \end{question}
}

\element{quizz}{
  \begin{question}{Q010}
    Variety ใน Big Data หมายถึงอะไร
    \begin{choices}
        \wrongchoice{ความเร็วของข้อมูล}
        \correctchoice{ชนิดและธรรมชาติของข้อมูลที่หลากหลาย}
        \wrongchoice{ปริมาณการจัดเก็บ}
        \wrongchoice{การใช้ cloud computing}
    \end{choices}
  \end{question}
}

\element{quizz}{
  \begin{question}{Q011}
    Velocity ของ Big Data หมายถึง
    \begin{choices}
        \correctchoice{ความเร็วในการสร้างข้อมูล}
        \wrongchoice{ความถูกต้องของข้อมูล}
        \wrongchoice{ความหลากหลายของไฟล์}
        \wrongchoice{ความแตกต่างระหว่างข้อมูลเชิงคุณภาพกับเชิงปริมาณ}
    \end{choices}
  \end{question}
}

\element{quizz}{
  \begin{question}{Q012}
    ไฟล์ Excel, CSV, Word, JSON จะถูกจัดในหมวดใด
    \begin{choices}
        \wrongchoice{Media}
        \correctchoice{Docs}
        \wrongchoice{Data Storage}
        \wrongchoice{Social Media}
    \end{choices}
  \end{question}
}

\element{quizz}{
  \begin{question}{Q013}
    Media vs Social Media ต่างกันอย่างไร
    \begin{choices}
        \correctchoice{Media = สื่อดิจิทัล, ส่วน Social Media = เครือข่ายออนไลน์}
        \wrongchoice{Media ต้องอยู่ในฐานข้อมูล SQL}
        \wrongchoice{Media เป็นข้อมูลเชิงตัวเลข, Social Media เป็นข้อมูลตาราง}
        \wrongchoice{ไม่มีความต่าง}
    \end{choices}
  \end{question}
}

\element{quizz}{
  \begin{question}{Q014}
    ข้อใดเป็นตัวอย่าง Sensor Data
    \begin{choices}
        \wrongchoice{Metadata ของรูปถ่าย}
        \correctchoice{ข้อมูลจาก Smart electric meters}
        \wrongchoice{Clickstream logs}
        \wrongchoice{Document repository}
    \end{choices}
  \end{question}
}

\element{quizz}{
  \begin{question}{Q015}
    ข้อใดคือการเรียงลำดับชั้นของ DIKW ที่ถูกต้อง
    \begin{choices}
        \wrongchoice{Wisdom -> Knowledge -> Information -> Data}
        \correctchoice{Data -> Information -> Knowledge -> Wisdom}
        \wrongchoice{Knowledge -> Wisdom -> Data -> Information}
        \wrongchoice{Information -> Data -> Wisdom -> Knowledge}
    \end{choices}
  \end{question}
}

\element{quizz}{
  \begin{question}{Q016}
    ถ้าเรามี “คะแนนสอบนักเรียน” แต่ยังไม่มีการตีความ ถือว่าอยู่ในระดับใดของ DIKW
    \begin{choices}
        \correctchoice{Data}
        \wrongchoice{Information}
        \wrongchoice{Knowledge}
        \wrongchoice{Wisdom}
    \end{choices}
  \end{question}
}

\element{quizz}{
  \begin{question}{Q017}
    ถ้าเรานำคะแนนสอบไปจัดเรียง ทำตาราง วิเคราะห์เป็นค่าเฉลี่ย ถือว่าอยู่ในระดับใดของ DIKW
    \begin{choices}
        \wrongchoice{Data}
        \correctchoice{Information}
        \wrongchoice{Knowledge}
        \wrongchoice{Wisdom}
    \end{choices}
  \end{question}
}

\element{quizz}{
  \begin{question}{Q018}
    ถ้าเรานำผลค่าเฉลี่ยไปอธิบายว่านักเรียนกลุ่มหนึ่งมีผลการเรียนอ่อนลงเพราะไม่ได้เตรียมสอบ ถือว่าอยู่ในระดับใดของ DIKW
    \begin{choices}
        \wrongchoice{Data}
        \wrongchoice{Information}
        \correctchoice{Knowledge}
        \wrongchoice{Wisdom}
    \end{choices}
  \end{question}
}

\element{quizz}{
  \begin{question}{Q019}
    ถ้าเรานำความรู้นี้ไปออกนโยบายปรับวิธีสอนเพื่อพัฒนาผลการเรียน ถือว่าอยู่ระดับใดของ DIKW
    \begin{choices}
        \wrongchoice{Data}
        \wrongchoice{Information}
        \wrongchoice{Knowledge}
        \correctchoice{Wisdom}
    \end{choices}
  \end{question}
}

% \element{quizz}{
%   \begin{question}{Q020}
%     Bruce Henderson (1975) เสนอแนวคิดเรื่องใด
%     \begin{choices}
%         \wrongchoice{Value Chain}
%         \correctchoice{Scale curves, Experience curves}
%         \wrongchoice{Five Forces Model}
%         \wrongchoice{Balanced Scorecard}
%     \end{choices}
%   \end{question}
% }

% \element{quizz}{
%   \begin{question}{Q021}
%     แนวคิด Scale curve และ Experience curve สื่อถึงอะไร
%     \begin{choices}
%         \correctchoice{การผลิตเพิ่มขึ้นทำให้ต้นทุนต่อหน่วยลดลง}
%         \wrongchoice{การผลิตมากขึ้นทำให้ต้นทุนสูงขึ้น}
%         \wrongchoice{ต้นทุนไม่เปลี่ยนแปลงตามปริมาณ}
%         \wrongchoice{การลงทุนไม่เกี่ยวกับต้นทุน}
%     \end{choices}
%   \end{question}
% }

% \element{quizz}{
%   \begin{question}{Q022}
%     Michael Porter (1980) เน้นแนวคิดใด
%     \begin{choices}
%         \correctchoice{Value Chain}
%         \wrongchoice{Balanced Scorecard}
%         \wrongchoice{Stakeholder Theory}
%         \wrongchoice{Game Theory}
%     \end{choices}
%   \end{question}
% }

% \element{quizz}{
%   \begin{question}{Q023}
%     ข้อใดคือความแตกต่างหลักของแนวคิด Henderson (1975) และ Porter (1980)
%     \begin{choices}
%         \correctchoice{Henderson = เน้นปริมาณ/ประสบการณ์กับต้นทุน, Porter = เน้นห่วงโซ่คุณค่าและโครงสร้างธุรกิจ}
%         \wrongchoice{Henderson = เน้น Value Chain, Porter = เน้น Economies of Mass}
%         \wrongchoice{Henderson = เน้น Balanced Scorecard, Porter = เน้น Game Theory}
%         \wrongchoice{ไม่มีความต่าง}
%     \end{choices}
%   \end{question}
% }

\element{quizz}{
  \begin{question}{Q024}
    Economies of Scale หมายถึงอะไร
    \begin{choices}
        \wrongchoice{ต้นทุนต่อหน่วยเพิ่มขึ้นเมื่อผลิตมากขึ้น}
        \correctchoice{ข้อได้เปรียบที่บริษัทขนาดใหญ่มีจากการลดต้นทุนเมื่อปริมาณเพิ่ม}
        \wrongchoice{การได้รับเงินอุดหนุนจากรัฐบาล}
        \wrongchoice{การผลิตน้อยแต่คุณภาพสูง}
    \end{choices}
  \end{question}
}

\element{quizz}{
  \begin{question}{Q025}
    Economies of Scale มีผลต่อธุรกิจอย่างไร
    \begin{choices}
        \wrongchoice{ทำให้ต้นทุนต่อหน่วยสูงขึ้นเมื่อผลิตเพิ่ม}
        \correctchoice{ทำให้ต้นทุนต่อหน่วยลดลงเมื่อผลิตเพิ่ม}
        \wrongchoice{ไม่มีผลต่อการผลิต}
        \wrongchoice{ทำให้ต้นทุนเพิ่มคงที่}
    \end{choices}
  \end{question}
}

\element{quizz}{
  \begin{question}{Q026}
    ข้อใดคือ Internal Economies of Scale
    \begin{choices}
        \wrongchoice{ได้รับสิทธิประโยชน์ทางภาษีจากรัฐบาล}
        \correctchoice{การซื้อวัตถุดิบปริมาณมากทำให้ราคาถูกลง}
        \wrongchoice{การสนับสนุนจากสมาคมอุตสาหกรรม}
        \wrongchoice{การช่วยเหลือจากต่างประเทศ}
    \end{choices}
  \end{question}
}

\element{quizz}{
  \begin{question}{Q027}
    ข้อใดคือ External Economies of Scale
    \begin{choices}
        \wrongchoice{การซื้อวัตถุดิบจำนวนมากโดยบริษัท}
        \correctchoice{การได้รับ preferential treatment จากรัฐบาล}
        \wrongchoice{การจัดการสายการผลิตภายในบริษัท}
        \wrongchoice{การรวมแผนกเพื่อประหยัดค่าใช้จ่าย}
    \end{choices}
  \end{question}
}

% \element{quizz}{
%   \begin{question}{Q028}
%     Diseconomies of Scale หมายถึงอะไร
%     \begin{choices}
%         \wrongchoice{การที่ขนาดเล็กเกินไปทำให้ไม่คุ้มทุน}
%         \correctchoice{การที่บริษัทขยายใหญ่เกินไปจนประสิทธิภาพลดลงและต้นทุนสูงขึ้น}
%         \wrongchoice{การแข่งขันทำให้ราคาตกต่ำ}
%         \wrongchoice{การที่บริษัทไม่ได้รับการสนับสนุนจากรัฐบาล}
%     \end{choices}
%   \end{question}
% }

\element{quizz}{
  \begin{question}{Q029}
    Psychographic segmentation เน้นปัจจัยใด
    \begin{choices}
        \wrongchoice{อายุ เพศ รายได้}
        \correctchoice{บุคลิกภาพ ไลฟ์สไตล์ ความเชื่อ}
        \wrongchoice{ที่อยู่ ภูมิศาสตร์}
        \wrongchoice{ความถี่การซื้อสินค้า}
    \end{choices}
  \end{question}
}

\element{quizz}{
  \begin{question}{Q030}
    ข้อใดคือความหมายของ Artificial Intelligence (AI)
    \begin{choices}
        \wrongchoice{ความสามารถของระบบที่เรียนรู้ได้เองโดยไม่ต้องเขียนโปรแกรมชัดเจน}
        \wrongchoice{การเรียนรู้ที่อาศัยโครงข่ายประสาทเทียมหลายชั้น}
        \correctchoice{โปรแกรมที่สามารถเรียนรู้และรับรู้ได้}
        \wrongchoice{การวิเคราะห์ข้อมูลเพื่อจัดกลุ่ม (Clustering)}
    \end{choices}
  \end{question}
}

% \element{quizz}{
%   \begin{question}{Q031}
%     Machine Learning (ML) แตกต่างจาก AI ทั่วไปอย่างไร
%     \begin{choices}
%         \wrongchoice{Machine Learning คือการเขียนโปรแกรมให้ทำงานแบบกำหนดกฎตายตัว}
%         \correctchoice{Machine Learning คืออัลกอริทึมที่ปรับปรุงประสิทธิภาพได้เองเมื่อเจอข้อมูลมากขึ้น}
%         \wrongchoice{Machine Learning เป็นระบบที่คิดเชิงเหตุผลเหมือนมนุษย์}
%         \wrongchoice{Machine Learning เน้นใช้ Big Data อย่างเดียว}
%     \end{choices}
%   \end{question}
% }

\element{quizz}{
  \begin{question}{Q032}
    ข้อใด ไม่ใช่ ลักษณะของ Deep Learning (DL)
    \begin{choices}
        \wrongchoice{เป็นสาขาย่อยของ Machine Learning}
        \wrongchoice{ใช้โครงข่ายประสาทเทียมหลายชั้น}
        \wrongchoice{ต้องการข้อมูลจำนวนมากในการฝึก}
        \correctchoice{สามารถทำงานได้โดยไม่ต้องอาศัยข้อมูลใด ๆ}
    \end{choices}
  \end{question}
}

\element{quizz}{
  \begin{question}{Q033}
    ตาม กฎของมัวร์ (Moore’s Law) กล่าวไว้อย่างไร
    \begin{choices}
        \wrongchoice{จำนวนทรานซิสเตอร์จะลดลงครึ่งหนึ่งทุก 2 ปี}
        \correctchoice{จำนวนทรานซิสเตอร์บนไมโครชิปจะเพิ่มขึ้น 2 เท่าทุก 2 ปี}
        \wrongchoice{ความเร็วของคอมพิวเตอร์จะลดลงทุก 5 ปี}
        \wrongchoice{ราคาคอมพิวเตอร์จะคงที่ตลอดเวลา}
    \end{choices}
  \end{question}
}

\element{quizz}{
  \begin{question}{Q034}
    ข้อใดอธิบาย Supervised Learning ได้ถูกต้องที่สุด
    \begin{choices}
        \wrongchoice{ปล่อยให้โมเดลค้นหาโครงสร้างข้อมูลเองโดยไม่มีคำตอบ}
        \wrongchoice{โมเดลเรียนรู้จากรางวัลและการลงโทษจากสภาพแวดล้อม}
        \correctchoice{โมเดลเรียนรู้จากตัวอย่างที่ถูกต้อง}
        \wrongchoice{โมเดลสุ่มทดลองโดยไม่ต้องประเมินผล}
    \end{choices}
  \end{question}
}

\element{quizz}{
  \begin{question}{Q035}
    ลักษณะสำคัญของ Unsupervised Learning คืออะไร
    \begin{choices}
        \wrongchoice{โมเดลสุ่มทดลองโดยไม่ต้องประเมินผล}
        \wrongchoice{เรียนรู้จากรางวัลและโทษ}
        \correctchoice{ไม่มีคำตอบให้ และให้โมเดลค้นหาโครงสร้างข้อมูลเอง}
        \wrongchoice{ใช้กฎที่มนุษย์เขียนไว้ล้วน ๆ}
    \end{choices}
  \end{question}
}

\element{quizz}{
  \begin{question}{Q036}
    ลักษณะสำคัญของ Reinforcement Learning คืออะไร
    \begin{choices}
        \wrongchoice{โมเดลสุ่มทดลองโดยไม่ต้องประเมินผล}
        \correctchoice{เรียนรู้จากการให้รางวัลและการลงโทษ}
        \wrongchoice{จัดกลุ่มข้อมูลเองโดยไม่รู้คำตอบ}
        \wrongchoice{ใช้กฎที่เขียนไว้}
    \end{choices}
  \end{question}
}

\element{quizz}{
  \begin{question}{Q043}
    หลักการสำคัญของ K-means clustering คืออะไร
    \begin{choices}
        \correctchoice{หาค่าเฉลี่ย ของแต่ละกลุ่มและจัดข้อมูลไปยังกลุ่มที่ใกล้ที่สุด}
        \wrongchoice{ใช้กฎที่มนุษย์เขียนเพื่อจัดกลุ่ม}
        \wrongchoice{ทำนายค่าตัวเลขต่อเนื่อง}
        \wrongchoice{ใช้รางวัลและโทษเพื่อฝึก AI}
    \end{choices}
  \end{question}
}

\element{quizz}{
  \begin{question}{Q044}
    AI และคอมพิวเตอร์มีผลกระทบต่อการจ้างงานในลักษณะใดบ้าง
    \begin{choices}
        \wrongchoice{ใช้แทนแรงงานในบางอาชีพ}
        \wrongchoice{การเปลี่ยนแปลงจะเกิดขึ้นอย่างรวดเร็ว}
        \wrongchoice{ลดจำนวนแรงงาน}
        \lastchoices
        \correctchoice{ถูกทุกข้อ}
    \end{choices}
  \end{question}
}

\element{quizz}{
  \begin{question}{Q045}
    ผลกระทบต่อแรงงานจาก AI ควรถูกมองอย่างไร
    \begin{choices}
        \wrongchoice{เป็นผลบวกเสมอ}
        \wrongchoice{เป็นผลลบเสมอ}
        \correctchoice{อาจเป็นได้ทั้งบวกและลบ ขึ้นอยู่กับการจัดการ}
        \wrongchoice{ไม่มีผลกระทบจริง}
    \end{choices}
  \end{question}
}

% \element{quizz}{
%   \begin{question}{Q046}
%     ความแตกต่างที่ถูกต้องระหว่าง Automation กับ Augmentation คืออะไร
%     \begin{choices}
%         \correctchoice{Automation = แทนที่งานบางส่วนด้วย AI/หุ่นยนต์, Augmentation = ใช้ AI ผู้ช่วยเพื่อเพิ่มประสิทธิภาพงาน}
%         \wrongchoice{ทั้งสองหมายถึงการเลิกจ้างพนักงานทั้งหมด}
%         \wrongchoice{Automation = เพิ่มจำนวนพนักงาน, Augmentation = ลดคุณภาพงาน}
%         \wrongchoice{Automation = ใช้คนล้วน, Augmentation = ไม่ใช้คนเลย}
%     \end{choices}
%   \end{question}
% }

% \element{quizz}{
%   \begin{question}{Q047}
%     คำว่า Blue-Collar หมายถึงอะไร
%     \begin{choices}
%         \wrongchoice{งานที่เกี่ยวข้องกับการทำงานในสำนักงาน}
%         \wrongchoice{งานด้านวิชาชีพ เช่น ทนาย แพทย์}
%         \correctchoice{งานใช้แรงงานหรือทักษะ เช่น ช่างไฟฟ้า ช่างยนต์}
%         \wrongchoice{งานที่เกี่ยวข้องกับการบริหารจัดการ}
%     \end{choices}
%   \end{question}
% }

% \element{quizz}{
%   \begin{question}{Q048}
%     คำว่า White-Collar หมายถึงอะไร
%     \begin{choices}
%         \wrongchoice{งานภาคการผลิต}
%         \correctchoice{งานที่ทำในออฟฟิศ เช่น ทนาย แพทย์ หรือพนักงานธนาคาร}
%         \wrongchoice{งานใช้แรงงานก่อสร้าง}
%         \wrongchoice{งานที่ไม่ใช้คอมพิวเตอร์}
%     \end{choices}
%   \end{question}
% }

% \element{quizz}{
%   \begin{question}{Q049}
%     เหตุใดงาน White-Collar บางประเภทจึงมีความเสี่ยงสูงในการถูกแทนที่ด้วย AI
%     \begin{choices}
%         \wrongchoice{เพราะเป็นงานที่ใช้แรงงาน}
%         \correctchoice{เพราะสามารถทำงานซ้ำ ๆ ตามขั้นตอนที่กำหนดได้}
%         \wrongchoice{เพราะไม่ต้องใช้ทักษะการคิดวิเคราะห์}
%         \wrongchoice{เพราะไม่ต้องใช้คอมพิวเตอร์}
%     \end{choices}
%   \end{question}
% }

\element{quizz}{
  \begin{question}{Q050}
    ข้อใดต่อไปนี้จัดอยู่ในกลุ่ม AI Services
    \begin{choices}
        \wrongchoice{Netflix}
        \correctchoice{Suno.com}
        \wrongchoice{Canva.com}
        \wrongchoice{MS Offices}
    \end{choices}
  \end{question}
}

% \element{quizz}{
%   \begin{question}{Q051}
%     ตัวอย่างของ Integrated AI คือข้อใด
%     \begin{choices}
%         \wrongchoice{Runwayml.com}
%         \wrongchoice{Napkin.ai}
%         \correctchoice{Notion.com}
%         \wrongchoice{Lovable.dev}
%     \end{choices}
%   \end{question}
% }

% \element{quizz}{
%   \begin{question}{Q052}
%     ข้อใดคือแพลตฟอร์มที่เกี่ยวข้องกับ Generative AI ด้านเพลง ในกลุ่ม AI Services
%     \begin{choices}
%         \wrongchoice{ChatGPT.com}
%         \correctchoice{Suno.com}
%         \wrongchoice{Runwayml.com}
%         \wrongchoice{Lovable.dev}
%     \end{choices}
%   \end{question}
% }

% \element{quizz}{
%   \begin{question}{Q053}
%     ข้อใดต่อไปนี้ไม่ใช่ตัวอย่างของ AI as a service (AIaaS)
%     \begin{choices}
%         \wrongchoice{OpenAI}
%         \wrongchoice{Google AI}
%         \correctchoice{Canva}
%         \wrongchoice{Amazon Web Services}
%     \end{choices}
%   \end{question}
% }

\element{quizz}{
  \begin{question}{Q054}
    Agentic AI คืออะไร
    \begin{choices}
        \wrongchoice{ไม่สามารถใช้ฐานข้อมูลภายนอกได้}
        \correctchoice{ตัดสินใจเลือกเครื่องมือเพื่อแก้ปัญหาได้เอง}
        \wrongchoice{ใช้ได้เฉพาะงานแปลภาษา}
        \wrongchoice{ต้องเขียนโค้ดทุกขั้นตอน}
    \end{choices}
  \end{question}
}

\element{quizz}{
  \begin{question}{Q055}
    จุดประสงค์หลักของ RAG คืออะไร
    \begin{choices}
        \wrongchoice{ทำให้โมเดลสร้างเนื้อหาจากศูนย์}
        \correctchoice{ทำให้โมเดลเข้าถึงเอกสารที่กำหนดไว้}
        \wrongchoice{ทำให้โมเดลลดขนาด parameter ลง}
        \wrongchoice{ทำให้โมเดลเรียนรู้แบบ unsupervised}
    \end{choices}
  \end{question}
}

% \element{quizz}{
%   \begin{question}{Q056}
%     คำว่า “Summarize” จัดอยู่ในส่วนใดของ Prompt Structure
%     \begin{choices}
%         \correctchoice{Instruction}
%         \wrongchoice{Context}
%         \wrongchoice{Input}
%         \wrongchoice{Output}
%     \end{choices}
%   \end{question}
% }

% \element{quizz}{
%   \begin{question}{Q057}
%     ข้อความ “a presentation on recent advances in machine learning” จัดอยู่ในส่วนใด
%     \begin{choices}
%         \wrongchoice{Instruction}
%         \correctchoice{Context}
%         \wrongchoice{Input}
%         \wrongchoice{Output}
%     \end{choices}
%   \end{question}
% }

% \element{quizz}{
%   \begin{question}{Q058}
%     ข้อความ “Machine learning is a rapidly evolving field that continues to transform industries by enabling systems to learn from data” จัดอยู่ในส่วนใด
%     \begin{choices}
%         \wrongchoice{Instruction}
%         \wrongchoice{Context}
%         \correctchoice{Input}
%         \wrongchoice{Output}
%     \end{choices}
%   \end{question}
% }

% \element{quizz}{
%   \begin{question}{Q059}
%     ข้อความ “a summary of no more than 150 words, highlighting the key advances mentioned.” จัดอยู่ในส่วนใด
%     \begin{choices}
%         \wrongchoice{Instruction}
%         \wrongchoice{Context}
%         \wrongchoice{Input}
%         \correctchoice{Output}
%     \end{choices}
%   \end{question}
% }

% \element{quizz}{
%   \begin{question}{Q060}
%     “Translate this English sentence into French.” เป็นตัวอย่างของเทคนิคใด
%     \begin{choices}
%         \wrongchoice{Contextual Guidance}
%         \correctchoice{Task Specific}
%         \wrongchoice{Domain Expertise}
%         \wrongchoice{Framing.}
%     \end{choices}
%   \end{question}
% }

% \element{quizz}{
%   \begin{question}{Q061}
%     “Write a short paragraph on New York City, highlighting its iconic landmarks.” จัดเป็นเทคนิคใด
%     \begin{choices}
%         \correctchoice{Contextual Guidance}
%         \wrongchoice{Task Specific}
%         \wrongchoice{User Feedback Loop}
%         \wrongchoice{Bias Mitigation}
%     \end{choices}
%   \end{question}
% }

% \element{quizz}{
%   \begin{question}{Q062}
%     “Explain the causes, symptoms, and treatments of hypothyroidism, including the latest research and medical guidelines.” จัดอยู่ในเทคนิคใด
%     \begin{choices}
%         \correctchoice{Domain Expertise}
%         \wrongchoice{Task Specific}
%         \wrongchoice{Framing}
%         \wrongchoice{Contextual Guidance}
%     \end{choices}
%   \end{question}
% }

% \element{quizz}{
%   \begin{question}{Q063}
%     “Write a 100-word paragraph on leadership traits without favoring any gender.” เป็นตัวอย่างของเทคนิคใด
%     \begin{choices}
%         \wrongchoice{User Feedback Loop}
%         \correctchoice{Bias Mitigation}
%         \wrongchoice{Task Specific}
%         \wrongchoice{Contextual Guidance}
%     \end{choices}
%   \end{question}
% }

% \element{quizz}{
%   \begin{question}{Q064}
%     “Summarize the article in 100 words, focusing on main findings and recommendations.” เป็นเทคนิคใด
%     \begin{choices}
%         \wrongchoice{Framing}
%         \correctchoice{Task Specific}
%         \wrongchoice{Bias Mitigation}
%         \wrongchoice{Domain Expertise}
%     \end{choices}
%   \end{question}
% }

% \element{quizz}{
%   \begin{question}{Q065}
%     ข้อใดอธิบาย Zero Shot Prompting ได้ถูกต้อง
%     \begin{choices}
%         \wrongchoice{การให้ตัวอย่างหลายตัวก่อนให้โมเดลตอบ}
%         \correctchoice{การให้คำสั่งทั่วไปโดยไม่มีตัวอย่างประกอบ}
%         \wrongchoice{การใช้โมเดลใหม่เสมอทุกครั้ง}
%         \wrongchoice{การฝึกโมเดลด้วยข้อมูลใหม่}
%     \end{choices}
%   \end{question}
% }

% \element{quizz}{
%   \begin{question}{Q066}
%     ตัวอย่าง “Translate ‘Hello’ to Thai” จัดอยู่ในรูปแบบใด
%     \begin{choices}
%         \wrongchoice{Few Shot}
%         \correctchoice{Zero Shot}
%         \wrongchoice{Domain Expertise}
%         \wrongchoice{Contextual Guidance}
%     \end{choices}
%   \end{question}
% }

% \element{quizz}{
%   \begin{question}{Q067}
%     ข้อใดเป็นตัวอย่างของ Few Shot Prompting
%     \begin{choices}
%         \wrongchoice{Translate “Hello” to Thai}
%         \correctchoice{Hello: สวัสดีเจ้า / Good bye: ไปก่อนเน้อเจ้า / Goodnight: …}
%         \wrongchoice{Write a short essay on AI}
%         \wrongchoice{Summarize the article in 100 words}
%     \end{choices}
%   \end{question}
% }

% \element{quizz}{
%   \begin{question}{Q068}
%     จุดเด่นของ Few Shot Prompting คืออะไร
%     \begin{choices}
%         \wrongchoice{ทำงานได้โดยไม่ต้องใช้ข้อมูลเลย}
%         \correctchoice{ให้โมเดลเรียนรู้จากตัวอย่างเพียงไม่กี่ตัวอย่าง}
%         \wrongchoice{ใช้ข้อมูลจำนวนมากในการฝึก}
%         \wrongchoice{ทำให้โมเดลไม่ต้องใช้ Context}
%     \end{choices}
%   \end{question}
% }

% \element{quizz}{
%   \begin{question}{Q069}
%     “ทําตัวเปนเชฟอาหารระดับโลก เดียวฉันจะบรรยายอาหารทีกิน แล้วคุณคอมเมนต์เกียวกับเมนู นันให้หน่อย” เป็นตัวอย่างของอะไร
%     \begin{choices}
%         \correctchoice{Persona Pattern}
%         \wrongchoice{Audience Persona Pattern}
%         \wrongchoice{Template Pattern}
%         \wrongchoice{Interview Pattern}
%     \end{choices}
%   \end{question}
% }

% \element{quizz}{
%   \begin{question}{Q070}
%     อธิบายการทํางานของห่วงโซ่อุปทานในร้านขายของชําของสหรัฐฯ ให้ฉันฟง สมมติว่าฉันคือ เจงกิสข่าน เป็นตัวอย่างของอะไร
%     \begin{choices}
%         \wrongchoice{Persona Pattern}
%         \correctchoice{Audience Persona Pattern}
%         \wrongchoice{Template Pattern}
%         \wrongchoice{Interview Pattern}
%     \end{choices}
%   \end{question}
% }

% \element{quizz}{
%   \begin{question}{Q071}
%     สร้างโปรแกรมออกกําลังกายตามเทมเพลตนี: ชือท่า, จํานวนครัง x เซ็ต, กลุ่มกล้ามเนือทีใช้, ระดับ ความยาก 1-5, หมายเหตุเรืองท่าทาง เป็นตัวอย่างของอะไร
%     \begin{choices}
%         \wrongchoice{Persona Pattern}
%         \wrongchoice{Audience Persona Pattern}
%         \correctchoice{Template Pattern}
%         \wrongchoice{Interview Pattern}
%     \end{choices}
%   \end{question}
% }

% \element{quizz}{
%   \begin{question}{Q072}
%     หากร้านขายปากกากล่องละ 5 ดอลลาร์ พร้อมโปรซื้อ 1 แถม 1 แล้วซื้อ 10 กล่อง ควรใช้ Pattern ใดในการแก้โจทย์
%     \begin{choices}
%         \correctchoice{Chain-of-Thought Pattern}
%         \wrongchoice{Few Shot Prompting}
%         \wrongchoice{Domain Expertise}
%         \wrongchoice{Bias Mitigation}
%     \end{choices}
%   \end{question}
% }

% \element{quizz}{
%   \begin{question}{Q073}
%     ตัวอย่างคำถาม “How to get a good grade” ใน Tree of Thought ใช้เพื่ออะไร
%     \begin{choices}
%         \wrongchoice{สาธิตการคำนวณเชิงตัวเลข}
%         \correctchoice{สาธิตการวิเคราะห์แบบหลายมุมมอง}
%         \wrongchoice{สาธิตการใช้ Zero-Shot Prompting}
%         \wrongchoice{สาธิตการลด Bias}
%     \end{choices}
%   \end{question}
% }

\element{quizz}{
  \begin{question}{Q074}
    เหตุผลที่ทำให้ IoT มีความสามารถในการประมวลผลและวิเคราะห์มากขึ้นคืออะไร
    \begin{choices}
        \correctchoice{การมีการเชื่อมต่อเซนเซอร์และอุปกรณ์เข้ากับอินเทอร์เน็ต}
        \wrongchoice{การตัดการเชื่อมต่อจากระบบเครือข่าย}
        \wrongchoice{การใช้พลังงานจากมนุษย์แทนเครื่องจักร}
        \wrongchoice{การลดจำนวนอุปกรณ์ที่เชื่อมต่อ}
    \end{choices}
  \end{question}
}

\element{quizz}{
  \begin{question}{Q075}
    IoT Gateway ทำหน้าที่อะไร
    \begin{choices}
        \wrongchoice{เป็นที่เก็บข้อมูลระยะยาว}
        \correctchoice{เป็นตัวเชื่อมต่ออุปกรณ์ IoT เข้าสู่เครือข่ายและคลาวด์}
        \wrongchoice{เป็นหน้าจอให้ผู้ใช้ควบคุม}
        \wrongchoice{เป็นอุปกรณ์ตรวจจับอุณหภูมิ}
    \end{choices}
  \end{question}
}

\element{quizz}{
  \begin{question}{Q076}
    Cloud ในระบบ IoT มีไว้ทำอะไร
    \begin{choices}
        \wrongchoice{แจกสัญญาณ Wi-Fi}
        \correctchoice{วิเคราะห์และจัดเก็บข้อมูลจากอุปกรณ์}
        \wrongchoice{ตรวจจับควันโดยตรง}
        \wrongchoice{เป็นที่เก็บข้อมูลระยะยาว}
    \end{choices}
  \end{question}
}

\element{quizz}{
  \begin{question}{Q077}
    “Internet of Things with IP address” หมายถึงอะไร
    \begin{choices}
        \correctchoice{อุปกรณ์ที่เชื่อมต่ออินเทอร์เน็ตได้เองโดยไม่ต้องพึ่งคน}
        \wrongchoice{อุปกรณ์ที่ต้องมีคนกดปุ่มทุกครั้ง}
        \wrongchoice{ระบบที่ใช้ได้เฉพาะคอมพิวเตอร์}
        \wrongchoice{อุปกรณ์ที่ไม่มีเซนเซอร์}
    \end{choices}
  \end{question}
}

\element{quizz}{
  \begin{question}{Q078}
    IoT without IP address หมายถึงอะไร
    \begin{choices}
        \wrongchoice{ระบบที่ไม่ใช้เซนเซอร์เลย}
        \correctchoice{ระบบที่มีเซนเซอร์ แต่ไม่เชื่อมต่อโดยตรงกับอินเทอร์เน็ต}
        \wrongchoice{ระบบที่มีเฉพาะมือถือและแท็บเล็ต}
        \wrongchoice{เครือข่ายที่ใช้เฉพาะ Wi-Fi}
    \end{choices}
  \end{question}
}

% \element{quizz}{
%   \begin{question}{Q079}
%     ชั้น Perception Layer มีหน้าที่หลักคืออะไร
%     \begin{choices}
%         \wrongchoice{วิเคราะห์ข้อมูลจำนวนมาก}
%         \correctchoice{ตรวจจับและเก็บข้อมูลจากสิ่งแวดล้อมด้วยเซนเซอร์}
%         \wrongchoice{จัดการเรื่องธุรกิจและความเป็นส่วนตัว}
%         \wrongchoice{ส่งต่อข้อมูลผ่านเครือข่าย}
%     \end{choices}
%   \end{question}
% }

% \element{quizz}{
%   \begin{question}{Q080}
%     ชั้น Transport Layer ทำหน้าที่ใด
%     \begin{choices}
%         \correctchoice{ส่งข้อมูลจากเซนเซอร์ไปยังชั้นอื่น ๆ ผ่านเครือข่าย}
%         \wrongchoice{ให้บริการเฉพาะเจาะจงกับผู้ใช้}
%         \wrongchoice{จัดการข้อมูลจำนวนมากบนคลาวด์}
%         \wrongchoice{ควบคุมโมเดลธุรกิจและการใช้งานระบบ}
%     \end{choices}
%   \end{question}
% }

% \element{quizz}{
%   \begin{question}{Q081}
%     Processing Layer ทำอะไร
%     \begin{choices}
%         \wrongchoice{เก็บข้อมูลบนอุปกรณ์ IoT โดยตรง}
%         \correctchoice{วิเคราะห์ ประมวลผล และจัดเก็บข้อมูลจำนวนมหาศาล}
%         \wrongchoice{ส่งข้อมูลต่อด้วยสัญญาณไร้สาย}
%         \wrongchoice{แสดงผลลัพธ์ให้ผู้ใช้}
%     \end{choices}
%   \end{question}
% }

% \element{quizz}{
%   \begin{question}{Q082}
%     ชั้น Application Layer มีหน้าที่ใด
%     \begin{choices}
%         \wrongchoice{ตรวจจับสภาพแวดล้อม}
%         \correctchoice{ให้บริการแอปพลิเคชันเฉพาะสำหรับผู้ใช้}
%         \wrongchoice{ควบคุมเครือข่ายการเชื่อมต่อ}
%         \wrongchoice{เก็บข้อมูลระยะยาว}
%     \end{choices}
%   \end{question}
% }

% \element{quizz}{
%   \begin{question}{Q083}
%     ชั้น Business Layer เกี่ยวข้องกับเรื่องใด
%     \begin{choices}
%         \correctchoice{โมเดลธุรกิจ ความเป็นส่วนตัว และการจัดการระบบโดยรวม}
%         \wrongchoice{เฉพาะการเชื่อมต่อ Bluetooth}
%         \wrongchoice{การตรวจจับข้อมูลสภาพอากาศ}
%         \wrongchoice{การประมวลผล Big Data โดยตรง}
%     \end{choices}
%   \end{question}
% }

% \element{quizz}{
%   \begin{question}{Q084}
%     ถ้าเป็นเครือข่ายเล็ก ๆ ในสำนักงานทั่วไปที่มีอุปกรณ์ไม่กี่ร้อยเครื่อง ควรใช้แบบใด
%     \begin{choices}
%         \correctchoice{IPv4}
%         \wrongchoice{IPv6}
%         \wrongchoice{ทั้งคู่ใช้ไม่ได้}
%         \wrongchoice{ต้องใช้เครือข่ายพิเศษเท่านั้น}
%     \end{choices}
%   \end{question}
% }

% \element{quizz}{
%   \begin{question}{Q085}
%     ถ้าต้องรองรับอุปกรณ์ IoT หลายพันล้านชิ้นทั่วโลก ควรใช้แบบใด
%     \begin{choices}
%         \wrongchoice{IPv4}
%         \correctchoice{IPv6}
%         \wrongchoice{ทั้งคู่ใช้ไม่ได้}
%         \wrongchoice{ต้องใช้เครือข่ายพิเศษเท่านั้น}
%     \end{choices}
%   \end{question}
% }

\element{quizz}{
  \begin{question}{Q086}
    ถ้าต้องการตรวจจับระดับแอลกอฮอล์ในลมหายใจ ควรใช้เซนเซอร์ใด
    \begin{choices}
        \wrongchoice{Gas Sensor}
        \correctchoice{Alcohol Sensor}
        \wrongchoice{PIR Sensor}
        \wrongchoice{Ultrasonic Sensor}
    \end{choices}
  \end{question}
}

\element{quizz}{
  \begin{question}{Q087}
    เซนเซอร์ใดใช้ตรวจสอบทิศทางและการหมุน
    \begin{choices}
        \correctchoice{Gyroscope Sensor}
        \wrongchoice{Proximity Sensor}
        \wrongchoice{Rain Sensor}
        \wrongchoice{Alcohol Sensor}
    \end{choices}
  \end{question}
}

\element{quizz}{
  \begin{question}{Q088}
    การเพิ่มนวัตกรรม (A) ส่งผลอย่างไรในโมเดล Solow model
    \begin{choices}
        \correctchoice{ทำให้เส้น Output (Y) ขยับขึ้น}
        \wrongchoice{ทำให้ทุน (K) ลดลง}
        \wrongchoice{ทำให้แรงงาน (eL) หายไป}
        \wrongchoice{ทำให้ค่าเสื่อมราคามากขึ้น}
    \end{choices}
  \end{question}
}

\element{quizz}{
  \begin{question}{Q089}
    คำว่า “Cloud technology” เกี่ยวข้องกับข้อใดมากที่สุด
    \begin{choices}
        \wrongchoice{ฮาร์ดดิสก์ในเครื่องคอมพิวเตอร์}
        \correctchoice{เครือข่ายหรืออินเทอร์เน็ต}
        \wrongchoice{โปรแกรมที่ติดตั้งในเครื่อง}
        \wrongchoice{แฟลชไดรฟ์ที่เสียบคอมพิวเตอร์}
    \end{choices}
  \end{question}
}

\element{quizz}{
  \begin{question}{Q090}
    Cloud ให้บริการผ่านเครือข่ายแบบใดได้บ้าง
    \begin{choices}
        \wrongchoice{ได้เฉพาะอินเทอร์เน็ตบ้าน}
        \correctchoice{ทั้งเครือข่ายสาธารณะและเครือข่ายส่วนตัว เช่น WAN, LAN, VPN}
        \wrongchoice{ได้เฉพาะบลูทูธ}
        \wrongchoice{ต้องต่อสายตรงระหว่างเครื่องเท่านั้น}
    \end{choices}
  \end{question}
}

% \element{quizz}{
%   \begin{question}{Q091}
%     ข้อใดเป็นตัวอย่าง application ที่มักรันบนคลาวด์
%     \begin{choices}
%         \wrongchoice{โปรแกรมเครื่องคิดเลขออฟไลน์}
%         \correctchoice{การส่งอีเมลและประชุมออนไลน์}
%         \wrongchoice{เกมออฟไลน์บนเครื่องเดียว}
%         \wrongchoice{โปรแกรมจดโน้ตแบบไม่เชื่อมเน็ต}
%     \end{choices}
%   \end{question}
% }

\element{quizz}{
  \begin{question}{Q092}
    ข้อใดคือความหมายของ Public Cloud
    \begin{choices}
        \wrongchoice{ใช้เฉพาะในองค์กรเดียว}
        \correctchoice{มีพื้นที่มากและขยายได้ง่าย เหมาะกับงานพัฒนาและทำงานร่วมกัน}
        \wrongchoice{ใช้หลายองค์กรร่วมกันเฉพาะในอุตสาหกรรมเดียว}
        \wrongchoice{ผสมผสานระหว่างระบบออฟไลน์กับออนไลน์}
    \end{choices}
  \end{question}
}

\element{quizz}{
  \begin{question}{Q093}
    Private Cloud เหมาะสมกับสถานการณ์ใดมากที่สุด
    \begin{choices}
        \wrongchoice{เกมออนไลน์ขนาดใหญ่}
        \correctchoice{ธุรกิจที่ต้องการความปลอดภัยสูง}
        \wrongchoice{การแชร์ข้อมูลระหว่างหลายองค์กรในอุตสาหกรรมเดียวกัน}
        \wrongchoice{นักเรียนที่อยากใช้พื้นที่ฟรี}
    \end{choices}
  \end{question}
}

\element{quizz}{
  \begin{question}{Q094}
    Hybrid Cloud มีจุดเด่นคืออะไร
    \begin{choices}
        \correctchoice{ผสมการใช้ Public และ Private Cloud เข้าด้วยกัน}
        \wrongchoice{ใช้เฉพาะกับการประชุมออนไลน์}
        \wrongchoice{เก็บข้อมูลได้เฉพาะในฮาร์ดดิสก์ภายในเครื่อง}
        \wrongchoice{ใช้สำหรับส่งอีเมลเท่านั้น}
    \end{choices}
  \end{question}
}

\element{quizz}{
  \begin{question}{Q095}
    Community Cloud หมายถึงข้อใด
    \begin{choices}
        \wrongchoice{ใช้แค่ภายในบ้าน}
        \correctchoice{แพลตฟอร์มที่หลายองค์กรในอุตสาหกรรมเดียวกันใช้ร่วมกัน}
        \wrongchoice{การเก็บข้อมูลบนมือถือ}
        \wrongchoice{ใช้พื้นที่จาก USB Drive}
    \end{choices}
  \end{question}
}

\element{quizz}{
  \begin{question}{Q096}
    Infrastructure as a Service (IaaS) คือบริการแบบใด
    \begin{choices}
        \correctchoice{ให้ผู้ใช้เช่าใช้ฮาร์ดแวร์และระบบพื้นฐานแบบตามต้องการ}
        \wrongchoice{ให้ผู้ใช้เฉพาะซอฟต์แวร์สำเร็จรูป}
        \wrongchoice{ให้ผู้ใช้ซื้อเครื่องเซิร์ฟเวอร์เอง}
        \wrongchoice{ให้ผู้ใช้เชื่อมต่อ Wi-Fi ฟรี}
    \end{choices}
  \end{question}
}

\element{quizz}{
  \begin{question}{Q097}
    Platform as a Service (PaaS) ให้บริการอะไรเป็นหลัก
    \begin{choices}
        \wrongchoice{ซอฟต์แวร์สำเร็จรูปพร้อมใช้งานทันที}
        \correctchoice{เครื่องมือสำหรับพัฒนาและรันแอปพลิเคชัน}
        \wrongchoice{เครื่องคอมพิวเตอร์จริงที่ต้องซื้อเอง}
        \wrongchoice{การเชื่อมต่ออินเทอร์เน็ตฟรี}
    \end{choices}
  \end{question}
}

\element{quizz}{
  \begin{question}{Q098}
    Software as a Service (SaaS) คืออะไร
    \begin{choices}
        \wrongchoice{การเช่าเครื่องคอมพิวเตอร์จริงจากผู้ให้บริการ}
        \correctchoice{การใช้ซอฟต์แวร์เป็นบริการผ่านอินเทอร์เน็ต}
        \wrongchoice{การซื้อโปรแกรมแล้วติดตั้งลงเครื่อง}
        \wrongchoice{การพัฒนาแพลตฟอร์มเองทั้งหมด}
    \end{choices}
  \end{question}
}

\element{quizz}{
  \begin{question}{Q099}
    บริการแบบ (Infrastructure as a Service) IaaS มักถูกใช้งานโดยใคร
    \begin{choices}
        \wrongchoice{ผู้ใช้งานทั่วไป}
        \wrongchoice{นักพัฒนาซอฟต์แวร์}
        \correctchoice{ผู้บรืหาร}
        \wrongchoice{นักเรียนมัธยม}
    \end{choices}
  \end{question}
}

\element{quizz}{
  \begin{question}{Q100}
    ตัวอย่างบริการใดตรงกับ SaaS
    \begin{choices}
        \correctchoice{Gmail, Trello, Office 365}
        \wrongchoice{AWS, Stackscale}
        \wrongchoice{Heroku, Cloud Foundry}
        \wrongchoice{VMware}
    \end{choices}
  \end{question}
}

\element{quizz}{
  \begin{question}{Q101}
    จุดเด่นสำคัญของ Edge Computing คืออะไร
    \begin{choices}
        \correctchoice{ช่วยลดเวลาในการตอบสนอง}
        \wrongchoice{ทำให้การประมวลผลช้าลง}
        \wrongchoice{ใช้ได้เฉพาะในเกม}
        \wrongchoice{ต้องมีอินเทอร์เน็ตตลอดเวลา เพื่อติดต่อภายนอก}
    \end{choices}
  \end{question}
}

\element{quizz}{
  \begin{question}{Q102}
    คำว่า Smart City ตามนิยามของ UK Parliament หมายถึงอะไร
    \begin{choices}
        \wrongchoice{เมืองที่มีตึกสูงและระบบขนส่งสะดวก}
        \correctchoice{เมืองที่ใช้เทคโนโลยีและข้อมูลมาช่วยพัฒนาเมือง}
        \wrongchoice{เมืองที่มีประชากรหนาแน่นและเศรษฐกิจดี}
        \wrongchoice{เมืองที่มีมรดกทางวัฒนธรรมเยอะ}
    \end{choices}
  \end{question}
}

\element{quizz}{
  \begin{question}{Q103}
    Automation ช่วยเมืองอัจฉริยะในด้านใด
    \begin{choices}
        \correctchoice{ทำให้ระบบสามารถวิเคราะห์ข้อมูลจากเซนเซอร์และสั่งการได้เอง}
        \wrongchoice{ทำให้เมืองมีระบบไฟถนนเพิ่มขึ้น}
        \wrongchoice{ทำให้เมืองสร้างอาคารสูงขึ้น}
        \wrongchoice{ทำให้ทุกคนต้องควบคุมเองทั้งหมด}
    \end{choices}
  \end{question}
}

\element{quizz}{
  \begin{question}{Q104}
    Sensors และ IoT (Internet of Things) ช่วยเมืองอัจฉริยะอย่างไร
    \begin{choices}
        \correctchoice{เก็บข้อมูลสิ่งแวดล้อม การจราจร คุณภาพอากาศ และเชื่อมต่ออุปกรณ์ต่าง ๆ}
        \wrongchoice{ทำให้ประชาชนใช้โซเชียลมีเดียได้มากขึ้น}
        \wrongchoice{สร้างความบันเทิง}
        \wrongchoice{ใช้แทนการวางผังเมืองทั้งหมด}
    \end{choices}
  \end{question}
}

\element{quizz}{
  \begin{question}{Q105}
    Smart Infrastructure มีบทบาทอย่างไรในเมืองอัจฉริยะ
    \begin{choices}
        \correctchoice{ระบบโครงสร้างพื้นฐานที่ช่วยสนับสนุนบริการดิจิทัล}
        \wrongchoice{การเลือกตั้งผู้ว่าราชการ}
        \wrongchoice{การใช้สมาร์ทโฟน}
        \wrongchoice{การจัดงานเทศกาล}
    \end{choices}
  \end{question}
}

\element{quizz}{
  \begin{question}{Q106}
    Smart Citizens ในเมืองอัจฉริยะหมายถึงอะไร
    \begin{choices}
        \correctchoice{พลเมืองที่มีความรู้ เท่าทันเทคโนโลยี และมีส่วนร่วมในการพัฒนาเมือง}
        \wrongchoice{ประชาชนที่ใช้สมาร์ทโฟนทุกคน}
        \wrongchoice{ประชาชนที่อยู่ในเมืองใหญ่}
        \wrongchoice{ประชาชนที่ทำงานด้านไอที}
    \end{choices}
  \end{question}
}

\element{quizz}{
  \begin{question}{Q107}
    Digital Infrastructure ใน Smart Infrastructure มีอะไรบ้าง
    \begin{choices}
        \wrongchoice{รถไฟฟ้า, ถนน, เขื่อน}
        \correctchoice{เซ็นเซอร์, IoT, เครือข่าย}
        \wrongchoice{โรงเรียน, โรงพยาบาล, ตลาด}
        \wrongchoice{แม่น้ำและภูเขา}
    \end{choices}
  \end{question}
}

\element{quizz}{
  \begin{question}{Q108}
    Physical Infrastructure ที่เป็นส่วนหนึ่งของ Smart Infrastructure ได้แก่ข้อใด
    \begin{choices}
        \correctchoice{พลังงาน น้ำ การขนส่ง ขยะ และโทรคมนาคม}
        \wrongchoice{เซนเซอร์ ปัญญาประดิษฐ์ คลาวด์ และบิ๊กดาต้า}
        \wrongchoice{สื่อสังคมออนไลน์ และแอปมือถือ}
        \wrongchoice{บล็อกเชน และโทเคน}
    \end{choices}
  \end{question}
}

% \element{quizz}{
%   \begin{question}{Q109}
%     โครงสร้างพื้นฐานแบบ Semi-intelligent มีลักษณะอย่างไร
%     \begin{choices}
%         \wrongchoice{เก็บข้อมูลและตัดสินใจเองได้}
%         \correctchoice{เก็บข้อมูลแต่ไม่สามารถตัดสินใจจากข้อมูลได้}
%         \wrongchoice{ไม่เก็บข้อมูลใด ๆ}
%         \wrongchoice{ตัดสินใจอัตโนมัติโดยไม่ต้องใช้มนุษย์}
%     \end{choices}
%   \end{question}
% }

\element{quizz}{
  \begin{question}{Q110}
    ตัวอย่างของ Semi-intelligent infrastructure คืออะไร
    \begin{choices}
        \wrongchoice{ระบบจราจรที่แนะนำเส้นทางให้ผู้ขับขี่}
        \correctchoice{แผนที่ที่บันทึกข้อมูลมลพิษหรือปริมาณรถในเมือง}
        \wrongchoice{ระบบไฟถนนที่ปรับแสงอัตโนมัติ}
        \wrongchoice{ระบบอาคารอัจฉริยะที่ควบคุมพลังงานเอง}
    \end{choices}
  \end{question}
}

% \element{quizz}{
%   \begin{question}{Q111}
%     Intelligent infrastructure ทำหน้าที่อย่างไร
%     \begin{choices}
%         \correctchoice{เก็บข้อมูลและนำเสนอเพื่อช่วยมนุษย์ตัดสินใจ}
%         \wrongchoice{ไม่เก็บข้อมูลใด ๆ}
%         \wrongchoice{ตัดสินใจได้เองโดยไม่ต้องมีมนุษย์}
%         \wrongchoice{ใช้ข้อมูลเพียงเพื่อการเก็บสถิติเท่านั้น}
%     \end{choices}
%   \end{question}
% }

% \element{quizz}{
%   \begin{question}{Q112}
%     Intelligent infrastructure ทำหน้าที่อย่างไร
%     \begin{choices}
%         \correctchoice{เก็บข้อมูลและนำเสนอเพื่อช่วยมนุษย์ตัดสินใจเอง}
%         \wrongchoice{ไม่เก็บข้อมูลใด ๆ}
%         \wrongchoice{ตัดสินใจได้เองโดยไม่ต้องมีมนุษย์}
%         \wrongchoice{ใช้ข้อมูลเพียงเพื่อการเก็บสถิติเท่านั้น}
%     \end{choices}
%   \end{question}
% }

% \element{quizz}{
%   \begin{question}{Q113}
%     ตัวอย่างของ Intelligent infrastructure คืออะไร
%     \begin{choices}
%         \wrongchoice{ระบบไฟถนนที่เปิดปิดเองตามเวลา}
%         \correctchoice{ระบบขนส่งที่ตรวจจับการจราจรติดขัดและแจ้งผู้ขับขี่}
%         \wrongchoice{แผนที่ที่บอกเฉพาะพื้นที่สีเขียว}
%         \wrongchoice{ระบบประปาที่เก็บข้อมูลการใช้น้ำแต่ไม่ปรับเปลี่ยนอะไร}
%     \end{choices}
%   \end{question}
% }

\element{quizz}{
  \begin{question}{Q114}
    Smart infrastructure แตกต่างจากประเภทอื่นตรงไหน
    \begin{choices}
        \wrongchoice{เก็บข้อมูลแต่ไม่ตัดสินใจ เพื่อให้มนุษย์ตัดสินใจเอง}
        \wrongchoice{ต้องใช้คนตัดสินใจเสมอ}
        \correctchoice{ตัดสินใจและปรับตัวได้เองแบบอัตโนมัติ โดยไม่ต้องมีมนุษย์}
        \wrongchoice{ใช้เทคโนโลยีแต่ไม่สามารถเปลี่ยนแปลงตามสภาพจริง}
    \end{choices}
  \end{question}
}

\element{quizz}{
  \begin{question}{Q115}
    ตามแนวคิด Smart City ของไทย มีกี่ด้าน ที่ใช้พัฒนาเมืองอัจฉริยะ
    \begin{choices}
        \wrongchoice{2 ด้าน}
        \wrongchoice{3 ด้าน}
        \correctchoice{7 ด้าน}
        \wrongchoice{11 ด้าน}
    \end{choices}
  \end{question}
}

\element{quizz}{
  \begin{question}{Q116}
    Smart Environment มีเป้าหมายหลักคืออะไร
    \begin{choices}
        \wrongchoice{การสร้างตึกสูงในเมือง}
        \correctchoice{การรักษาสมดุลธรรมชาติและใช้ทรัพยากรอย่างยั่งยืน}
        \wrongchoice{การสร้างสนามบินใหม่}
        \wrongchoice{การเพิ่มปริมาณยานพาหนะในเมือง}
    \end{choices}
  \end{question}
}

\element{quizz}{
  \begin{question}{Q117}
    ข้อใดคือ ตัวอย่างของ Smart Energy
    \begin{choices}
        \wrongchoice{การสร้างสวนสาธารณะ}
        \wrongchoice{ระบบขนส่งสาธารณะ}
        \correctchoice{พลังงานหมุนเวียน และระบบ Smart Grid}
        \wrongchoice{การออกแบบเมืองเพื่อการท่องเที่ยว}
    \end{choices}
  \end{question}
}

\element{quizz}{
  \begin{question}{Q118}
    Smart Mobility เน้นในเรื่องใดมากที่สุด
    \begin{choices}
        \correctchoice{การเดินทางและระบบขนส่งที่สะดวก ปลอดภัย และเป็นมิตรต่อสิ่งแวดล้อม}
        \wrongchoice{การสร้างตลาดและแหล่งชุมชนใหม่}
        \wrongchoice{การใช้พลังงานไฟฟ้าในครัวเรือน}
        \wrongchoice{การสร้างสวนดอกไม้}
    \end{choices}
  \end{question}
}

\element{quizz}{
  \begin{question}{Q119}
    ข้อใด ไม่ใช่ หลักการของ Smart Living
    \begin{choices}
        \wrongchoice{Smart Health (สุขภาพและสุขอนามัยของคน)}
        \wrongchoice{Public Safety and Security (ความปลอดภัยจากอาชญากรรมและภัยพิบัติ)}
        \wrongchoice{Smart Built Environment (สิ่งแวดล้อมการอยู่อาศัยที่ชาญฉลาด)}
        \correctchoice{Smart Agriculture (เกษตรอัจฉริยะ)}
    \end{choices}
  \end{question}
}

% \element{quizz}{
%   \begin{question}{Q120}
%     Smart Environment มุ่งเน้นด้านใดมากที่สุด
%     \begin{choices}
%         \wrongchoice{การพัฒนาเทคโนโลยีสารสนเทศ}
%         \correctchoice{การรักษาสมดุลธรรมชาติและการใช้ทรัพยากรอย่างยั่งยืน}
%         \wrongchoice{การสร้างระบบไฟฟ้าพลังงานใหม่}
%         \wrongchoice{การพัฒนาเขตเศรษฐกิจพิเศษ}
%     \end{choices}
%   \end{question}
% }

% \element{quizz}{
%   \begin{question}{Q121}
%     Smart People มุ่งพัฒนาเรื่องใด
%     \begin{choices}
%         \correctchoice{การใช้เทคโนโลยีดิจิทัลเพื่อเรียนรู้ตลอดชีวิตและลดความเหลื่อมล้ำ}
%         \wrongchoice{การสร้างพลังงานไฟฟ้าใหม่}
%         \wrongchoice{การจัดการระบบขนส่ง}
%         \wrongchoice{การพัฒนาภาคเกษตรกรรม}
%     \end{choices}
%   \end{question}
% }

\element{quizz}{
  \begin{question}{Q122}
    Smart Governance เกี่ยวข้องกับอะไร
    \begin{choices}
        \correctchoice{การบริหารภาครัฐที่โปร่งใส}
        \wrongchoice{การสร้างพลังงานหมุนเวียน}
        \wrongchoice{การพัฒนาเมืองท่องเที่ยว}
        \wrongchoice{การสร้างระบบไฟฟ้าแรงสูง}
    \end{choices}
  \end{question}
}

\element{quizz}{
  \begin{question}{Q123}
    Smart Economy มีเป้าหมายเพื่ออะไร
    \begin{choices}
        \wrongchoice{การสร้างสวนสาธารณะในเมือง}
        \correctchoice{การส่งเสริมธุรกิจใหม่ ๆ และระบบเศรษฐกิจดิจิทัล}
        \wrongchoice{การพัฒนาพลังงานไฟฟ้า}
        \wrongchoice{การพัฒนาภาคเกษตรกรรมแบบดั้งเดิม}
    \end{choices}
  \end{question}
}

% \element{quizz}{
%   \begin{question}{Q124}
%     ข้อใดเป็นส่วนหนึ่งของ Vision & Goals
%     \begin{choices}
%         \wrongchoice{การใช้พลังงานหมุนเวียน}
%         \correctchoice{การกำหนดเขต Smart City และเป้าหมายของเมือง}
%         \wrongchoice{การทำระบบขนส่งอัจฉริยะ}
%         \wrongchoice{การเก็บภาษีในเมือง}
%     \end{choices}
%   \end{question}
% }


% \element{quizz}{
%   \begin{question}{Q126}
%     City Data & Security ต้องมีอย่างน้อย 3 เรื่อง อะไรบ้าง
%     \begin{choices}
%         \correctchoice{City Data Catalogue, City Data Exchange, City Governance}
%         \wrongchoice{Smart Economy, Smart Mobility, Smart Energy}
%         \wrongchoice{E-Government, E-Banking, E-Commerce}
%         \wrongchoice{Big Data, Blockchain, AI}
%     \end{choices}
%   \end{question}
% }

% \element{quizz}{
%   \begin{question}{Q127}
%     ใน 7 Dimensions of Smart City Solutions สิ่งที่บังคับว่าต้องมีคืออะไร
%     \begin{choices}
%         \correctchoice{Smart Economy}
%         \wrongchoice{Smart Environment}
%         \wrongchoice{Smart People}
%         \wrongchoice{Smart Governance}
%     \end{choices}
%   \end{question}
% }

% \element{quizz}{
%   \begin{question}{Q128}
%     Management for Sustainability มุ่งเน้นเรื่องใด
%     \begin{choices}
%         \wrongchoice{การจัดเก็บภาษีท้องถิ่น}
%         \correctchoice{การมีแผนและกลไกเพื่อความยั่งยืน}
%         \wrongchoice{การขยายเมืองออกไปสู่ชานเมือง}
%         \wrongchoice{การสร้างสนามกีฬาในเมือง}
%     \end{choices}
%   \end{question}
% }



\element{final}{
  \begin{question}{QF001}
    ข้อใดอธิบายคุณลักษณะ Veracity ของ Big Data ได้ดีที่สุด
    \begin{choices}
        \correctchoice{ข้อมูลมีความถูกต้องและเชื่อถือได้}
        \wrongchoice{ข้อมูลมีปริมาณมหาศาล}
        \wrongchoice{ข้อมูลมีการเปลี่ยนแปลงอย่างรวดเร็ว}
        \wrongchoice{ข้อมูลมีความหลากหลายของรูปแบบ}
    \end{choices}
  \end{question}
}

\element{final}{
  \begin{question}{QF002}
    การวิเคราะห์ข้อมูล Big Data มีประโยชน์อย่างไรต่อองค์กรทางธุรกิจมากที่สุด
    \begin{choices}
        \correctchoice{ทำให้ได้เปรียบทางการแข่งขันและตัดสินใจแม่นยำขึ้น}
        \wrongchoice{ลดพื้นที่การจัดเก็บข้อมูล}
        \wrongchoice{ทำให้พนักงานทำงานน้อยลง}
        \wrongchoice{สามารถลบข้อมูลเก่าทิ้งได้ทั้งหมด}
    \end{choices}
  \end{question}
}

\element{final}{
  \begin{question}{QF003}
    ข้อความใดอธิบาย Value ใน Big Data ได้ถูกต้อง
    \begin{choices}
        \wrongchoice{ข้อมูลต้องมีปริมาณมากพอ}
        \correctchoice{ข้อมูลที่สามารถนำมาสกัดเป็นความรู้และสร้างมูลค่าได้}
        \wrongchoice{ข้อมูลที่ไหลเข้าสู่ระบบอย่างรวดเร็ว}
        \wrongchoice{ข้อมูลที่มีการเรียงลำดับอย่างเป็นระเบียบ}
    \end{choices}
  \end{question}
}

\element{final}{
  \begin{question}{QF004}
    เครื่องมือใดที่มักใช้ในการรับและจัดเก็บข้อมูลระดับ Big Data ที่ไม่มีโครงสร้างจำเพาะ
    \begin{choices}
        \wrongchoice{Microsoft Word}
        \correctchoice{NoSQL หรือ เครือข่ายแบบ Hadoop}
        \wrongchoice{Relational Database}
        \wrongchoice{Notepad}
    \end{choices}
  \end{question}
}

\element{final}{
  \begin{question}{QF005}
    ประโยชน์หลักของ Data visualization คืออะไร
    \begin{choices}
        \wrongchoice{ทำให้เก็บข้อมูลได้เร็วขึ้น}
        \correctchoice{นำเสนอข้อมูลที่ซับซ้อนให้เข้าใจง่ายผ่านภาพกราฟิก}
        \wrongchoice{ลดข้อผิดพลาดของข้อมูลต้นทาง}
        \wrongchoice{ทดแทนการใช้ AI}
    \end{choices}
  \end{question}
}

% \element{final}{
%   \begin{question}{QF006}
%     ความแตกต่างของ Data และ Information ตามแนวคิด DIKW คืออะไร
%     \begin{choices}
%         \correctchoice{Data เป็นข้อมูลดิบ, Information คือข้อมูลที่ผ่านการประมวลผลแล้ว}
%         \wrongchoice{Information คือข้อมูลดิบ, Data คือความรู้ที่นำไปใช้}
%         \wrongchoice{ไม่มีข้อแตกต่าง}
%         \wrongchoice{Data ต้องเป็นตัวเลขเท่านั้น ส่วน Information เป็นข้อความ}
%     \end{choices}
%   \end{question}
% }

\element{final}{
  \begin{question}{QF007}
    Wisdom ในลำดับชั้น DIKW หมายถึงข้อใด
    \begin{choices}
        \wrongchoice{การแปลงประโยคเป็นตาราง}
        \correctchoice{การนำความรู้ไปประยุกต์ใช้ในการตัดสินใจและสร้างกลยุทธ์}
        \wrongchoice{ข้อมูลดิบที่ได้จากเซนเซอร์บนตัวเครื่อง}
        \wrongchoice{การตรวจสอบความถูกต้องของข้อมูลทั้งหมด}
    \end{choices}
  \end{question}
}

\element{final}{
  \begin{question}{QF008}
    การเก็บสถิติยอดขายรายวันเป็นตัวเลขเก็บไว้เฉย ๆ เป็นตัวอย่างของระดับใดใน DIKW
    \begin{choices}
        \wrongchoice{Wisdom}
        \wrongchoice{Knowledge}
        \wrongchoice{Information}
        \correctchoice{Data}
    \end{choices}
  \end{question}
}

% \element{final}{
%   \begin{question}{QF009}
%     กลยุทธ์การเป็นผู้นำด้านต้นทุน (Cost Leadership) ตามแนวคิดของ Michael Porter เน้นประเด็นใด
%     \begin{choices}
%         \wrongchoice{การหาสินค้าที่โดดเด่นไม่เหมือนใคร}
%         \correctchoice{การบริหารต้นทุนการผลิตและดำเนินการให้ต่ำกว่าคู่แข่ง}
%         \wrongchoice{เน้นขายสินค้าราคาสูงที่สุดโดยไม่แคร์ตลาด}
%         \wrongchoice{การกำหนดกลุ่มเป้าหมายลูกค้าคนชั้นสูงเป็นหลัก}
%     \end{choices}
%   \end{question}
% }

% \element{final}{
%   \begin{question}{QF010}
%     ในโมเดล Value Chain ของ Porter กิจกรรมใดจัดเป็น Primary Activities
%     \begin{choices}
%         \wrongchoice{การบริหารทรัยากรบุคคล (HR Management)}
%         \wrongchoice{การพัฒนาเทคโนโลยี (Technology Development)}
%         \correctchoice{การผลิต (Operations) และ โลจิสติกส์ขาออก (Outbound Logistics)}
%         \wrongchoice{การจัดหาจัดซื้อวัตถุดิบ (Procurement)}
%     \end{choices}
%   \end{question}
% }

\element{final}{
  \begin{question}{QF011}
    ปัจจัยหลักสำคัญอะไรที่ส่งเสริมให้เกิด Economies of Scale ขึ้นในการผลิตรายใหญ่
    \begin{choices}
        \wrongchoice{ค่าเช่าที่ดินที่สูงขึ้นทุกวัน}
        \correctchoice{ต้นทุนคงที่ (Fixed Cost) ปันส่วนเฉลี่ยไปในปริมาณที่ผลิตเพิ่มขึ้น}
        \wrongchoice{การจ้างพนักงานพาร์ทไทม์มากขึ้นเพื่อลดรายจ่าย}
        \wrongchoice{การลดวัตถุดิบลงแต่ปรับราคาสินค้าขึ้น}
    \end{choices}
  \end{question}
}

\element{final}{
  \begin{question}{QF012}
    หลักการของ Learning Curve (Experience Curve) สะท้อนให้เห็นรูปผลลัพธ์ใดทางธุรกิจ
    \begin{choices}
        \wrongchoice{พนักงานจะขอขึ้นเงินเดือนเสมอเมื่อเวลาผ่านไป}
        \correctchoice{องค์กรมีประสิทธิภาพสูงขึ้น ต้นทุนต่อหน่วยลดลงตามการสั่งสมประสบการณ์}
        \wrongchoice{ผู้บริโภคจะเบื่อหน่ายสินค้าเร็วขึ้นเสมอ}
        \wrongchoice{ผู้บริหารไม่ต้องควบคุมงานด้วยตนเอง}
    \end{choices}
  \end{question}
}

\element{final}{
  \begin{question}{QF013}
    ปรากฏการณ์ Diseconomies of Scale มักสะท้อนปัญหาที่เป็นเกิดจากสาเหตุใด
    \begin{choices}
        \correctchoice{องค์กรขยายขนาดจนเกิดความซับซ้อนและการบริหารขาดประสิทธิภาพ}
        \wrongchoice{บริษัทใช้เครื่องจักรล้ำสมัยเกินไปทำให้เกิดต้นทุนสูง}
        \wrongchoice{การมีผู้บริหารที่เก่งเกินไป}
        \wrongchoice{ขายดีจนสินค้าขาดสต็อกอย่างรวดเร็ว}
    \end{choices}
  \end{question}
}

% \element{final}{
%   \begin{question}{QF014}
%     ข้อใดเป็นคำนิยาม Machine Learning (ML) ที่เห็นภาพได้กระชับและถูกต้องที่สุด
%     \begin{choices}
%         \wrongchoice{ระบบที่ทำงานเฉพาะตามคำสั่งเงื่อนไข IF/THEN ขั้นพื้นฐาน}
%         \correctchoice{ระบบที่สามารถหาแพทเทิร์นเรียนรู้และปรับปรุงตนเองจากชุดข้อมูลโดยไม่ต้องเขียนโค้ดบรรทัดต่อบรรทัด}
%         \wrongchoice{ระบบวิเคราะห์ข้อมูลสูตรตาราง Excel แบบตายตัวซ้ำๆ}
%         \wrongchoice{เป็นเพียงเครื่องจักรแขนกลโรงงานที่มีสมองสั่งงาน}
%     \end{choices}
%   \end{question}
% }

\element{final}{
  \begin{question}{QF015}
    กระบวนการใดใน ML ที่โมเดลต้องเรียนรู้จากการลองผิดลองถูกเพื่อบรรลุผลและรับรางวัล
    \begin{choices}
        \wrongchoice{Supervised Learning}
        \wrongchoice{Unsupervised Learning}
        \correctchoice{Reinforcement Learning}
        \wrongchoice{Transfer Learning}
    \end{choices}
  \end{question}
}

\element{final}{
  \begin{question}{QF016}
    เทคนิค K-means ถูกใช้งานกับการแก้ปัญหาประเภทใด
    \begin{choices}
        \correctchoice{การจัดกลุ่มชุดข้อมูลที่ไม่มีป้ายกำกับ}
        \wrongchoice{หาความสัมพันธ์พยากรณ์ราคาทองคำ}
        \wrongchoice{สั่งการระบบขับขี่รถยนต์ไร้คนขับ}
        \wrongchoice{ประมวลผลแปลเสียงคนและสุนัข}
    \end{choices}
  \end{question}
}

\element{final}{
  \begin{question}{QF017}
    จุดที่ Deep Learning โดดเด่นกว่า ML ปกติทั่วไปในด้านโครงสร้างคืออะไร
    \begin{choices}
        \correctchoice{มีโครงข่ายประสาทเทียมหลายชั้นลึก ช่วยวิเคราะห์แพทเทิร์นความซับซ้อนสูงได้ดี}
        \wrongchoice{เป็นโมเดลที่แทบไม่ต้องใช้ข้อมูลใด ๆ}
        \wrongchoice{ฝึกได้ง่ายที่สุดและใช้เวลาน้อยมากเพียงไม่กี่วินาทีบนคอมธรรมดา}
        \wrongchoice{ใช้งานง่ายแค่โยนข้อมูลภาพดิบ ก็สามารถล้างไวรัสได้}
    \end{choices}
  \end{question}
}

\element{final}{
  \begin{question}{QF018}
    ในบริบทของ AI, คำว่า NLP (Natural Language Processing) มุ่งเน้นไปที่ทักษะใดมากที่สุด
    \begin{choices}
        \wrongchoice{การวิเคราะห์และดึงเป้าหมายจากภาพถ่ายดาวเทียมทางภูมิศาสตร์}
        \correctchoice{การเข้าใจ วิเคราะห์ และโต้ตอบกับภาษามนุษย์}
        \wrongchoice{การตั้งจุดสมดุลการขับขี่พวงมาลัยรถ}
        \wrongchoice{แยกแยะสีของเสื้อผ้าสัตว์เลี้ยงในบ้าน}
    \end{choices}
  \end{question}
}

\element{final}{
  \begin{question}{QF019}
    ระบบใดทำงานโดยใช้หลักปัญญาประดิษฐ์ด้าน Computer Vision
    \begin{choices}
        \wrongchoice{ระบบตรวจสอบไวยากรณ์ใน Word}
        \correctchoice{ระบบสแกนใบหน้าปลดล็อคมือถือ}
        \wrongchoice{แอปพลิเคชันพอดแคสต์}
        \wrongchoice{ระบบควบคุมอุณหภูมิโรงงานอัตโนมัติ}
    \end{choices}
  \end{question}
}

\element{final}{
  \begin{question}{QF020}
    หากมีเทคโนโลยีแบบ Augmentation AI ทำงานอยู่ มันจะมีหน้าที่ในข้อใดชัดที่สุด
    \begin{choices}
        \wrongchoice{ทำหน้าที่ตัดความรับผิดชอบของมนุษย์ออกไป}
        \correctchoice{ทำงานร่วมเพื่อเสริมทักษะ และเป็นผู้ช่วยช่วยให้มนุษย์ตัดสินใจเร็วขึ้นดีขึ้น}
        \wrongchoice{ยุติบทบาทการตัดสินใจของมนุษย์เมื่อถึงเวลาคับขัน}
        \wrongchoice{ทำงานเฉพาะเมื่ออยู่ในภาวะวิกฤตที่ต้องมีการควบคุมฮาร์ดแวร์}
    \end{choices}
  \end{question}
}

% \element{final}{
%   \begin{question}{QF021}
%     Generative AI แตกต่างจาก AI แบบวิเคราะห์แบบเดิมตรงข้อใด
%     \begin{choices}
%         \wrongchoice{ป้องกันการแพร่มัลแวร์}
%         \correctchoice{สามารถสร้างคอนเทนต์ใหม่เช่น ข้อความ, ภาพ, โค้ดที่คล้ายมนุษย์สร้างได้จากข้อมูลที่เรียนรู้}
%         \wrongchoice{สร้างข้อความลบข้อมูลที่ไม่ดีบนอินเตอร์เน็ต}
%         \wrongchoice{พยากรณ์ตัวเลขการเติบโตของเศรษฐกิจได้ถูกต้อง 100%}
%     \end{choices}
%   \end{question}
% }

\element{final}{
  \begin{question}{QF022}
    ใน AI ปรากฏการณ์คำว่า AI Hallucination ระบุถึงอาการแบบใดในฝั่งโมเดลภาษาขนาดใหญ่
    \begin{choices}
        \wrongchoice{AI ไม่สามารถโต้ตอบเมื่อเจอคำหยาบ}
        \correctchoice{AI ให้คำตอบที่ดูมีเหตุผลสวยงามแต่เนื้อหาเป็นข้อมูลเท็จ}
        \wrongchoice{เครื่อง Server AI เกิดอาการร้อนจัด}
        \wrongchoice{AI ขอเรียนรู้คำศัพท์ผู้ใช้บันทึกใส่ระบบตลอดเวลา}
    \end{choices}
  \end{question}
}

% \element{final}{
%   \begin{question}{QF023}
%     ตัวอย่างของ Generative AI ด้านการเจาะจงที่ใช้ในการช่วยเขียนเนื้อหาข้อความ (Text) คืออะไร
%     \begin{choices}
%         \wrongchoice{DALL-E}
%         \wrongchoice{Midjourney}
%         \correctchoice{ChatGPT}
%         \wrongchoice{Photoshop AI}
%     \end{choices}
%   \end{question}
% }

\element{final}{
  \begin{question}{QF024}
    เซนเซอร์หลักชนิดใดที่เกี่ยวข้องระบบสมาร์ทโฮมด้านการเปิดไฟเมื่อมีคนเดินผ่าน
    \begin{choices}
        \wrongchoice{Ultrasonic Sensor}
        \correctchoice{Motion Sensor}
        \wrongchoice{Alcohol Sensor}
        \wrongchoice{Gyroscope Sensor}
    \end{choices}
  \end{question}
}

\element{final}{
  \begin{question}{QF025}
    สถาปัตยกรรมโครงข่าย IoT ต้องมีสิ่งประกอบที่ตอบโจทย์หลายส่วน ข้อใดคือกระบวนการหลักที่ส่งเสริมเป็นลำดับ
    \begin{choices}
        \correctchoice{อุปกรณ์ฮาร์ดแวร์/เซนเซอร์, เครือข่าย, ระบบวิเคราะห์ข้อมูล, แอพพลิเคชั่นโต้ตอบกับผู้ใช้}
        \wrongchoice{เฉพาะจอภาพ ลำโพง และการเชื่อมต่อปลั๊กเพาเวอร์}
        \wrongchoice{โทรศัพท์ พาวเวอร์แบงค์ และเราเตอร์ขนาดจิ๋ว}
        \wrongchoice{คอมพิวเตอร์ระดับสูง ปลั๊กพ่วง บลูทูธเวอร์ชันล่าสุด}
    \end{choices}
  \end{question}
}

% \element{final}{
%   \begin{question}{QF026}
%     ระบบ IoT ในภาคการเกษตร (Smart Agriculture) มักเน้นการติดตั้งอุปกรณ์ใดเป็นฟังก์ชั่นหลัก
%     \begin{choices}
%         \wrongchoice{อุปกรณ์ GPS ติดปลอกคอสุนัข}
%         \correctchoice{เซนเซอร์วัดค่าความชื้นดินและปริมาณแสงแดดเพื่อจ่ายน้ำอัตโนมัติ}
%         \wrongchoice{สมาร์ทแบนด์สำหรับเกษตรกรเพื่อควบคุมหัวใจ}
%         \wrongchoice{แผงป้ายไฟแจ้งการเก็บเกี่ยวพืชในเมืองขนาดใหญ่}
%     \end{choices}
%   \end{question}
% }

\element{final}{
  \begin{question}{QF027}
    ข้อใด ไม่ใช่ ความท้าทายหลักของระบบการวางโครงข่าย IoT ในปัจจุบัน
    \begin{choices}
        \wrongchoice{การรักษาความปลอดภัยจากแฮกเกอร์}
        \correctchoice{การจัดหาอุปกรณ์ที่ต่ออินเทอร์เน็ตได้}
        \wrongchoice{ปัญหาความหน่วงสัญญาณ (Latency)}
        \wrongchoice{การจัดการพลังงานแบตเตอรี่ในอุปกรณ์ไร้สายขนาดเล็ก}
    \end{choices}
  \end{question}
}

% \element{final}{
%   \begin{question}{QF028}
%     มาตรฐานการเชื่อมต่อสัญญาณใดที่ทำงานได้ในระยะทางไกลแต่ประหยัดพลังงานสำหรับอุปกรณ์ IoT
%     \begin{choices}
%         \wrongchoice{Wi-Fi 6}
%         \correctchoice{LoRaWAN / NB-IoT}
%         \wrongchoice{HDMI Cable}
%         \wrongchoice{Direct USB Protocol}
%     \end{choices}
%   \end{question}
% }

\element{final}{
  \begin{question}{QF029}
    จุดมุ่งหมายสูงสุดของการใช้ Edge Computing คือช่วยบรรเทาภาระจุดใดเป็นหลักใน Cloud IoT
    \begin{choices}
        \wrongchoice{ปัญหาความจุปริมาณการจัดเก็บที่เต็มเร็ว}
        \correctchoice{ลดความหน่วง (Latency) ในการสื่อสารและการพึ่งพาเซิร์ฟเวอร์บนคลาวด์}
        \wrongchoice{ลดจำนวนการใช้เทอร์โมสตัทในบ้านลง}
        \wrongchoice{ป้องกันความจำเครื่องคอมพิวเตอร์หายแบบถาวร}
    \end{choices}
  \end{question}
}

\element{final}{
  \begin{question}{QF030}
    ทำไม IPv6 จึงสำคัญต่อระบบนิเวศของ IoT มหาศาลกว่า IPv4 เดิมมากมาย
    \begin{choices}
        \wrongchoice{เพราะโหลดรูปภาพเร็วกว่าในอินเทอร์เน็ต}
        \correctchoice{รองรับจำนวน IP Addresses มหาศาลเพียงพอต่อการเชื่อมต่อทุกอุปกรณ์}
        \wrongchoice{เพราะมีการออกแบบมาแก้ปัญหาเน็ตหลุดบ่อย}
        \wrongchoice{เนื่องจาก IPv6 ถูกผูกขาดลิขสิทธิ์ฟรี}
    \end{choices}
  \end{question}
}

\element{final}{
  \begin{question}{QF031}
    Actuator หรือเครื่องตอบสนองในแวดวง IoT มีลักษณะรูปแบบงานแบบใดตามทฤษฎี
    \begin{choices}
        \wrongchoice{รับค่าความเย็นของอากาศภายนอก}
        \correctchoice{เปลี่ยนแปลงสัญญาณคำสั่งกลับกลายมาเป็นพลังงานทางกล}
        \wrongchoice{ปรับเปลี่ยนรหัสไฟล์ให้คอมไพล์ทำงานได้}
        \wrongchoice{แสดงผลกราฟบนจอ}
    \end{choices}
  \end{question}
}

\element{final}{
  \begin{question}{QF032}
    ข้อใดคือแพลตฟอร์มที่เป็นลักษณะของ SaaS ที่ใช้ในงานประมวลผลเอกสารกลุ่มบริษัท
    \begin{choices}
        \wrongchoice{AWS EC2 Infrastructure}
        \correctchoice{Google Workspace, Office 365, Microsoft OneDrive}
        \wrongchoice{C++ Framework Console}
        \wrongchoice{Ubuntu OS Server edition}
    \end{choices}
  \end{question}
}

\element{final}{
  \begin{question}{QF033}
    บริการเช่าเฉพาะทรัพยากรพื้นฐาน เช่น คอมพิวเตอร์สเปคเปล่า ๆ เครือข่าย และพื้นที่สำหรับเก็บข้อมูล (Storage) นั้นตรงกับประเภทบริการคลาวด์ใด
    \begin{choices}
        \wrongchoice{SaaS}
        \wrongchoice{PaaS}
        \correctchoice{IaaS}
        \wrongchoice{DaaS}
    \end{choices}
  \end{question}
}

\element{final}{
  \begin{question}{QF034}
    ความเสี่ยงที่มักเป็นปัญหากังวลหลักขององค์กรเมื่อพิจารณาเลือกใช้ Public Cloud คือสิ่งใด
    \begin{choices}
        \wrongchoice{แบนวิดท์น้อยเกินไป}
        \wrongchoice{ต้องรอวิศวกรซ่อมบำรุง 3 เดือน}
        \correctchoice{ประเด็นความรอบคอบในความปลอดภัย}
        \wrongchoice{ผู้พัฒนาไม่เพิ่มพื้นที่การเก็บข้อมูลเลย}
    \end{choices}
  \end{question}
}

% \element{final}{
%   \begin{question}{QF035}
%     ในมุมมองสถาปัตยกรรม Serverless Computing มีข้อดีที่โดดเด่นในระบบคลาวด์คืออะไร
%     \begin{choices}
%         \wrongchoice{องค์กรลงทุนซื้อระบบไม่มีเซิร์ฟเวอร์}
%         \correctchoice{นักพัฒนาไม่ต้องวุ่นวายกับการจัดการเซิร์ฟเวอร์ โดยจะถูกคิดเงินเฉพาะส่วนโค้ดฟังก์ชั่นที่มีการรันเบื้องหลังเท่านั้น}
%         \wrongchoice{ทำให้ระบบทำงานได้แม้อินเตอร์เน็ตตัดขาดนานแรมเดือน}
%         \wrongchoice{สามารถจำกัดปริมาณโค้ดที่รันให้คงที่เสมอ}
%     \end{choices}
%   \end{question}
% }

\element{final}{
  \begin{question}{QF036}
    เป้าหมายองค์ประกอบด้าน Smart Mobility ในคอนเซ็ปต์ Smart City พุ่งเป้าหมายการแก้ปัญหาใดหลัก
    \begin{choices}
        \wrongchoice{ขยายตัวเมืองออกขอบปริมณฑล}
        \correctchoice{ผลักดันให้การจราจรราบรื่น}
        \wrongchoice{สนับสนุนการเดินเท้าทุกสถานที่}
        \wrongchoice{ยกเลิกระบบถนนทั้งหมด}
    \end{choices}
  \end{question}
}

% \element{final}{
%   \begin{question}{QF037}
%     โครงสร้างพื้นฐานประเภทใด ที่สามารถ วิเคราะห์ประมวลข้อมูล และปรับปรุงแก้ปัญหาระบบด้วยตนเอง อย่างแท้จริง
%     \begin{choices}
%         \wrongchoice{Physical infrastructure}
%         \wrongchoice{Semi-intelligent infrastructure}
%         \correctchoice{Intelligent / Smart infrastructure}
%         \wrongchoice{Traditional infrastructure}
%     \end{choices}
%   \end{question}
% }

\element{final}{
  \begin{question}{QF038}
    คำนิยาม Digital Twin ในแง่บริบทชุมชนเมือง คือเรื่องใด
    \begin{choices}
        \wrongchoice{ระบบสร้างตึกแฝดแบบกระจกโปร่งใส 2 ตึกเพื่อต้านลม เพื่อตรวจดูสถานการณ์จำลองปัญหาต่าง ๆ}
        \correctchoice{เมืองจำลองเชิงดิจิทัลที่ทำงานควบคู่กับการรับข้อมูลจากเซนเซอร์ในเมืองจริง}
        \wrongchoice{แอปพลิเคชันสอนภาษาและโปรแกรมมิ่ง}
        \wrongchoice{ใบรับรองอนุญาตสิทธิพลเมืองออนไลน์}
    \end{choices}
  \end{question}
}

\element{final}{
  \begin{question}{QF039}
    ปัจจัยตัวแปรหลักอะไรเป็นองค์ประกอบสำคัญต่อความสำเร็จของการทำ Smart City
    \begin{choices}
        \wrongchoice{พละกำลังของฮาร์ดแวร์คลาวด์ที่ทันสมัย}
        \correctchoice{ความร่วมแรงร่วมใจของหน่วยงานรัฐ ภาคเอกชน และประชาชนที่มีส่วนร่วมรับมือเทคโนโลยีไปด้วยกันอย่างยั่งยืน}
        \wrongchoice{การดึงบริษัทขนาดใหญ่นอกทวีปผู้เดียวมาทำโครงสร้าง}
        \wrongchoice{การวางงบที่ไม่จำกัดจากคลังกระทรวงมหาดไทย}
    \end{choices}
  \end{question}
}

% \element{final}{
%   \begin{question}{QF040}
%     การพัฒนา Smart Economy มุ่งหมายเพื่อยกระดับโครงสร้างใด
%     \begin{choices}
%         \wrongchoice{ทำให้ทุกฝ่ายกลับมาใช้วิธีซื้อและโอนมูลค่าผ่านธนบัตรเดิมเพื่อให้ตรวจสอบง่าย}
%         \correctchoice{สนับสนุนนวัตกรรมสร้างโมเดลธุรกิจใหม่ ส่งเสริม e-commerce เพื่อดึงจุดแข็งด้านการแข่งขันเทคโนโลยีดิจิทัล}
%         \wrongchoice{เร่งควบคุมกฎเกณฑ์ลดสวัสดิการของแรงงานสายเทค 50%}
%         \wrongchoice{เพิ่มโควต้าการประหยัดไฟร้านสะดวกซื้อในตอนกลางคืน}
%     \end{choices}
%   \end{question}
% }

% \element{final}{
%   \begin{question}{QF041}
%     ข้อใดกล่าวสอดคล้องถึงลักษณะองค์กรธุรกิจที่มีความรับผิดชอบ Management for Sustainability
%     \begin{choices}
%         \wrongchoice{ลดภาษีให้ระดับฝ่ายจัดการระดับสูง}
%         \correctchoice{มีกลไกตรวจสอบระบบการทำงาน ขยายพัฒนาความรู้ด้านเทคโนโลยีเป็นมิตรกับสิ่งแวดล้อมเพื่อเตรียมรับมืออนาคตระยะยาว}
%         \wrongchoice{มุ่งเน้นความสวยงามของตึกสูงสำนักงาน}
%         \wrongchoice{กระจายสาขาร้านอาหารอย่างเร่งด่วนที่สุดในสเกลอำเภอขนาดเล็ก}
%     \end{choices}
%   \end{question}
% }

% \element{final}{
%   \begin{question}{QF042}
%     ในโมเดล Solow-Swan ตัวแปรเทคโนโลยีหรือนวัตกรรม (A) มีบทบาทเช่นไรกับการผลิตมวลรวมในระยะยาวประการใด
%     \begin{choices}
%         \wrongchoice{เร่งเกิดค่าเสื่อมทุนเร็วกว่ากำหนด 3 เท่า}
%         \correctchoice{กระตุ้นเส้นสมรรถนะของปริมาณมวลรวมผลผลิต หรือ Output ให้ยกชูสูงขึ้นเมื่อเทียบในสัดส่วนความประหยัดเดิม}
%         \wrongchoice{ขัดขวางสมการความเจริญเติบโตกดดันให้ประเทศย้ายจุดดุลยภาพ}
%         \wrongchoice{ส่งผลให้ปริมาณแรงงานทิ้งระบบเศรษฐกิจโดยถาวร}
%     \end{choices}
%   \end{question}
% }

\element{final}{
  \begin{question}{QF043}
    AWS ให้บริการผ่านเครือข่ายแบบใด
    \begin{choices}
        \wrongchoice{เป็นการติดตั้งแอปจัดการความรู้ที่ราคาประหยัด}
        \correctchoice{มีโซลูชั่นรองรับระดับ IaaS และ PaaS}
        \wrongchoice{ออกแบบเฉพาะให้องค์กรเล็ก ๆ ที่ไม่ใช่อินเทอร์เน็ตทำงานเอกสาร}
        \wrongchoice{สร้างเครื่องจักรและเซนเซอร์ประกอบรถยนต์แบรนด์อเมซอน}
    \end{choices}
  \end{question}
}

% \element{final}{
%   \begin{question}{QF044}
%     เมืองที่มีความเป็นอัจฉริยะ (Smart City) ย่อมใส่ใจเรื่อง Smart Environment ด้วยเทคโนโลยีใดบ้างที่เป็นรูปธรรม
%     \begin{choices}
%         \wrongchoice{การจัดการระบบไฟถนนกระพริบตามจังหวะสัญญาณโทรศัพท์มือถือส่วนบุคคล}
%         \correctchoice{การติดตั้งระบบเซนเซอร์วัดค่าฝุ่นละอองและการจำแนกขยะให้เหมาะสมกับการจัดการทางทรัพยากรหมุนเวียน}
%         \wrongchoice{การออกแบบรถยนต์ส่วนตัวทุกคันให้มีถุงลม 12 ใบ}
%         \wrongchoice{ลดเวลาทำงานออฟฟิศลงและปิดไฟหน้าร้านชอปปิ้งทั้งหมด}
%     \end{choices}
%   \end{question}
% }


% new q uesiotns
\element{final}{
  \begin{question}{QF046}
    บริษัท Startup แห่งหนึ่งต้องการพัฒนาแอปพลิเคชันมือถือแบบรวดเร็ว โดยไม่ต้องการเสียเวลาตั้งค่าระบบปฏิบัติการ ฐานข้อมูล หรือเซิร์ฟเวอร์ด้วยตนเอง ควรเลือกใช้บริการคลาวด์รูปแบบใด
    \begin{choices}
        \wrongchoice{IaaS}
        \correctchoice{PaaS}
        \wrongchoice{SaaS}
        \wrongchoice{DaaS}
    \end{choices}
  \end{question}
}

\element{final}{
  \begin{question}{QF047}
    หากโรงเรียนต้องการระบบอีเมลและการจัดการเอกสารออนไลน์ให้บุคลากรทุกคนใช้งานได้ทันทีโดยไม่ต้องติดตั้งโปรแกรมหรือพัฒนาระบบเอง ควรเลือกใช้บริการใด
    \begin{choices}
        \wrongchoice{เช่า Server เปล่า ๆ}
        \wrongchoice{เช่าแพลตฟอร์มสำหรับเขียนโค้ด}
        \correctchoice{ใช้งาน Google Workspace หรือ Microsoft 365}
        \wrongchoice{ซื้อ Hardware Server มาตั้งที่โรงเรียน}
    \end{choices}
  \end{question}
}

\element{final}{
  \begin{question}{QF048}
    ธนาคารแห่งหนึ่งต้องการย้ายระบบคำนวณดอกเบี้ยที่มีความซับซ้อนและเขียนด้วยภาษาคอมพิวเตอร์แบบเฉพาะเจาะจง ขึ้นสู่คลาวด์ โดยธนาคารต้องการควบคุมตั้งแต่ระดับระบบปฏิบัติการ (OS) และหน่วยความจำ (RAM) อย่างเต็มที่ บริการใดตอบโจทย์ที่สุด
    \begin{choices}
        \wrongchoice{SaaS}
        \wrongchoice{PaaS}
        \correctchoice{IaaS}
        \wrongchoice{Serverless Computing}
    \end{choices}
  \end{question}
}

\element{final}{
  \begin{question}{QF049}
    นักพัฒนาคนหนึ่งใช้บริการ Heroku หรือ Google App Engine ในการนำโค้ดที่เขียนเสร็จแล้วขึ้นไปทำงานบนเซิร์ฟเวอร์ทันที โดยไม่ต้องดูแลเรื่องการอัปเดตระบบปฏิบัติการ นี่คือตัวอย่างของการใช้คลาวด์ระดับใด
    \begin{choices}
        \wrongchoice{IaaS}
        \correctchoice{PaaS}
        \wrongchoice{SaaS}
        \wrongchoice{FaaS}
    \end{choices}
  \end{question}
}

\element{final}{
  \begin{question}{QF050}
    บริษัทกราฟิกดีไซน์เช่าใช้งาน Adobe Creative Cloud รายเดือนเพื่อให้พนักงานทุกคนใช้งานโปรแกรม Photoshop และ Illustrator ได้ ทิศทางการลงทุนนี้จัดเก็บรูปแบบใด
    \begin{choices}
        \correctchoice{SaaS}
        \wrongchoice{PaaS}
        \wrongchoice{IaaS}
        \wrongchoice{STaaS}
    \end{choices}
  \end{question}
}

\element{final}{
  \begin{question}{QF051}
    ข้อใดเป็นข้อเสียหรือความท้าทายหลักที่องค์กรพบจากการใช้งาน SaaS
    \begin{choices}
        \wrongchoice{ต้องเสียค่าบำรุงรักษาเซิร์ฟเวอร์ด้วยตนเองสูงมาก}
        \wrongchoice{ต้องจ้างนักพัฒนาโปรแกรมเมอร์ขั้นสูงจำนวนมากมาเขียนโค้ด}
        \correctchoice{จำกัดการปรับแต่งให้เข้ากับกระบวนการทำงานแบบเฉพาะเจาะจงขององค์กร}
        \wrongchoice{ใช้เวลาในการติดตั้งฮาร์ดแวร์เบื้องต้นนานกว่าจะใช้งานได้}
    \end{choices}
  \end{question}
}

\element{final}{
  \begin{question}{QF052}
    แพลตฟอร์มอย่าง Salesforce หรือ ระบบจัดการลูกค้าสัมพันธ์ (CRM) ออนไลน์ทั่วไป ที่พนักงานฝ่ายขายสามารถล็อกอินเข้าไปกรอกข้อมูลและดูยอดขายได้เลย เป็นการให้บริการคลาวด์รูปแบบใด
    \begin{choices}
        \wrongchoice{IaaS}
        \wrongchoice{PaaS}
        \correctchoice{SaaS}
        \wrongchoice{On-Premises}
    \end{choices}
  \end{question}
}

\element{final}{
  \begin{question}{QF053}
    ถ้าบริษัทอีคอมเมิร์ซต้องการจัดการช่วงเทศกาลลดราคา ที่ผู้ใช้จะเพิ่มขึ้น 10 เท่าชั่วคราว จึงเช่า Virtual Machine (VM) เพิ่มเติมจาก AWS เพื่อรันระบบขององค์กร นี่คือประโยชน์เด่นของคลาวด์ระดับ IaaS ข้อใด
    \begin{choices}
        \wrongchoice{การไม่ต้องใช้โปรแกรมเมอร์}
        \correctchoice{ความยืดหยุ่นในการขยายและลดทรัพยากรได้อย่างรวดเร็ว}
        \wrongchoice{การได้โปรแกรมบัญชีฟรี}
        \wrongchoice{การลดเวลาการเขียนโค้ดลง}
    \end{choices}
  \end{question}
}

\element{final}{
  \begin{question}{QF054}
    กรณีใดที่องค์กร ไม่ควร เลือกใช้ SaaS
    \begin{choices}
        \wrongchoice{ต้องการใช้งานซอฟต์แวร์ทันทีภายในไม่กี่ชั่วโมง}
        \wrongchoice{มีบุคลากรไอทีน้อยและไม่ต้องการดูแลระบบเบื้องหลัง}
        \correctchoice{องค์กรมีกระบวนการทำงานที่ซับซ้อนมาก}
        \wrongchoice{ต้องการจ่ายเงินแบบเช่าใช้เป็นรายเดือนตามจำนวนผู้ใช้จริง}
    \end{choices}
  \end{question}
}

\element{final}{
  \begin{question}{QF055}
    บริการด้าน Cloud รูปแบบใดที่ทำหน้าที่เปรียบเสมือน ห้องครัว ที่เตรียมอุปกรณ์เตาอบ โต๊ะเตรียมอาหาร และเครื่องครัวไว้ให้เชฟเข้ามาทำอาหาร โดยไม่ต้องสร้างห้องครัวใหม่เอง
    \begin{choices}
        \wrongchoice{SaaS}
        \correctchoice{PaaS}
        \wrongchoice{IaaS}
        \wrongchoice{Public Cloud}
    \end{choices}
  \end{question}
}

\element{final}{
  \begin{question}{QF056}
    การเช่าพื้นที่เก็บข้อมูล และพลังประมวลผล บน Microsoft Azure เพื่อลงระบบปฏิบัติการ Linux เอง และติดตั้งฐานข้อมูลเอง ถือเป็นการใช้งานแบบใด
    \begin{choices}
        \correctchoice{IaaS}
        \wrongchoice{PaaS}
        \wrongchoice{SaaS}
        \wrongchoice{FaaS}
    \end{choices}
  \end{question}
}

\element{DLT}{
  \begin{question}{QF057}
    Distributed Ledger Technology (DLT) คืออะไร
    \begin{choices}
        \wrongchoice{ระบบเก็บข้อมูลแบบรวมศูนย์ที่ปลอดภัยที่สุด}
        \wrongchoice{ซอฟต์แวร์สำหรับจำลองการทำงานของเซิร์ฟเวอร์}
        \correctchoice{ระบบฐานข้อมูลที่แชร์และซิงโครไนซ์ข้อมูล}
        \wrongchoice{เครือข่ายโซเชียลมีเดียรูปแบบใหม่}
    \end{choices}
  \end{question}
}

\element{DLT}{
  \begin{question}{QF058}
    เทคโนโลยีใดเป็นที่รู้จักมากที่สุดซึ่งเป็นรูปแบบหนึ่งของ Distributed Ledger Technology (DLT)
    \begin{choices}
        \correctchoice{Blockchain}
        \wrongchoice{Cloud Storage}
        \wrongchoice{Big Data}
        \wrongchoice{IoT}
    \end{choices}
  \end{question}
}

\element{DLT}{
  \begin{question}{QF059}
    ความแตกต่างหลักระหว่าง Distributed Ledger Technology (DLT) กับฐานข้อมูลแบบดั้งเดิมคืออะไร
    \begin{choices}
        \wrongchoice{DLT ทำงานเร็วกว่าฐานข้อมูลแบบดั้งเดิมเสมอ}
        \correctchoice{DLT ไม่มีผู้ดูแลศูนย์กลางเพียงผู้เดียว}
        \wrongchoice{DLT จำเป็นต้องใช้บนโทรศัพท์มือถือ}
        \wrongchoice{DLT ใช้งานได้ฟรี}
    \end{choices}
  \end{question}
}

\element{DLT}{
  \begin{question}{QF060}
    ข้อใดคือความหมายของคำว่า Decentralization
    \begin{choices}
        \wrongchoice{การรวมศูนย์ข้อมูลไว้ที่เซิร์ฟเวอร์เดียว}
        \wrongchoice{การลบข้อมูลที่ไม่จำเป็นทิ้งโดยอัตโนมัติ}
        \correctchoice{การกระจายอำนาจและการควบคุมออกจากศูนย์กลาง}
        \wrongchoice{การเข้ารหัสข้อมูลด้วยรหัสผ่านเดียว}
    \end{choices}
  \end{question}
}

\element{DLT}{
  \begin{question}{QF061}
    Blockchain ถูกออกแบบมาให้มีคุณสมบัติหลักข้อใดที่โดดเด่น
    \begin{choices}
        \wrongchoice{ข้อมูลถูกลบอัตโนมัติเมื่อครบกำหนดเวลา}
        \correctchoice{ข้อมูลที่ถูกบันทึกแล้วไม่สามารถแก้ไขหรือลบได้ง่าย}
        \wrongchoice{ขนาดไฟล์มีขนาดเล็กตลอดเวลา}
        \wrongchoice{ความเร็วในการประมวลผลเร็วกว่าฐานข้อมูลทั่วไปหลายเท่า}
    \end{choices}
  \end{question}
}

\element{DLT}{
  \begin{question}{QF062}
    Block ในระบบ Blockchain ทำหน้าที่อะไร
    \begin{choices}
        \correctchoice{เก็บชุดของข้อมูลหรือธุรกรรมที่เกิดขึ้น}
        \wrongchoice{ประมวลผลกราฟิกของระบบ}
        \wrongchoice{เชื่อมต่ออินเทอร์เน็ตให้กับโหนด}
        \wrongchoice{ลบข้อมูลที่ซ้ำซ้อน}
    \end{choices}
  \end{question}
}

\element{DLT}{
  \begin{question}{QF063}
    สิ่งใดที่ใช้เชื่อมโยงแต่ละ Block ใน Blockchain เข้าด้วยกันแบบเป็นลำดับ
    \begin{choices}
        \wrongchoice{IP Address ของผู้สร้าง}
        \wrongchoice{ลายเซ็นดิจิทัลของผู้ดูแลระบบ}
        \correctchoice{Hash ของบล็อกก่อนหน้า}
        \wrongchoice{พาสเวิร์ดที่ใช้ร่วมกัน}
    \end{choices}
  \end{question}
}

\element{DLT}{
  \begin{question}{QF064}
    ฟังก์ชันแฮช (Hash Function) ที่ใช้ใน Blockchain มีลักษณะอย่างไร
    \begin{choices}
        \wrongchoice{แปลงข้อมูลกลับไปเป็นข้อความต้นฉบับได้เสมอ}
        \correctchoice{แปลงข้อมูลรับเข้าขนาดเท่าใดก็ได้ให้เป็นค่าที่มีความยาวคงที่}
        \wrongchoice{แปลงหรือลดขนาดไฟล์ข้อมูลภาพและเสียง}
        \wrongchoice{อนุญาตให้ทุกคนแก้ไขผลลัพธ์ได้}
    \end{choices}
  \end{question}
}

\element{DLT}{
  \begin{question}{QF065}
    หากมีการเปลี่ยนแปลงข้อมูลแม้เพียงเล็กน้อยใน Block ผลลัพธ์ของ Hash ของบล็อกนั้นจะเป็นอย่างไร
    \begin{choices}
        \wrongchoice{Hash จะเปลี่ยนไปเพียงเล็กน้อยตามสัดส่วน}
        \wrongchoice{Hash จะยังคงเหมือนเดิม}
        \wrongchoice{ระบบจะซ่อน Hash ทันทีเพื่อความปลอดภัย}
        \correctchoice{Hash จะเปลี่ยนไปอย่างสิ้นเชิง}
    \end{choices}
  \end{question}
}

\element{DLT}{
  \begin{question}{QF066}
    Consensus Algorithm ในเทคโนโลยี DLT มีไว้เพื่อเป้าหมายใด
    \begin{choices}
        \wrongchoice{เพื่อเข้ารหัสข้อมูลบัญชีผู้ใช้งานใหม่}
        \wrongchoice{เพื่อบีบอัดข้อมูลก่อนส่งผ่านอินเทอร์เน็ต}
        \correctchoice{เพื่อให้ทุกโหนดในเครือข่ายตกลงและยอมรับความถูกต้องของข้อมูลร่วมกัน}
        \wrongchoice{เพื่อป้องกันไวรัสคอมพิวเตอร์เข้าสู่ระบบ}
    \end{choices}
  \end{question}
}

\element{DLT}{
  \begin{question}{QF067}
    ข้อใดคือกลไกฉันทามติ (Consensus Mechanism) ดั้งเดิมที่ใช้ในเครือข่าย Bitcoin
    \begin{choices}
        \wrongchoice{Proof of Stake (PoS)}
        \correctchoice{Proof of Work (PoW)}
        \wrongchoice{Proof of Authority (PoA)}
        \wrongchoice{Delegated Proof of Stake (DPoS)}
    \end{choices}
  \end{question}
}

\element{DLT}{
  \begin{question}{QF068}
    กลไก Proof of Work (PoW) ต้องอาศัยสิ่งใดเป็นหลักในการแข่งขันกันเพื่อสร้างบล็อกใหม่
    \begin{choices}
        \correctchoice{พลังการประมวลผลของเครื่องคอมพิวเตอร์}
        \wrongchoice{จำนวนเหรียญที่ถือครอง}
        \wrongchoice{เครดิตสกอร์ของผู้ใช้งานจากธนาคาร}
        \wrongchoice{ความเร็วของการเชื่อมต่ออินเทอร์เน็ตของผู้ดูแลระบบ}
    \end{choices}
  \end{question}
}

\element{DLT}{
  \begin{question}{QF069}
    กลไก Proof of Stake (PoS) แตกต่างจาก Proof of Work (PoW) อย่างไร
    \begin{choices}
        \wrongchoice{ใช้ไฟฟ้าในการคำนวณมากกว่า PoW หลายเท่า}
        \wrongchoice{เป็นระบบที่มีหน่วยงานกลางของรัฐควบคุมโดยสมบูรณ์}
        \correctchoice{ใช้จำนวนเหรียญที่วางค้ำประกัน (Stake) เพื่อสิทธิในการสร้างบล็อกแทนการแข่งขันด้านพลังประมวลผล}
        \wrongchoice{อนุญาตให้ทุกคนแก้ไขบล็อกเก่าได้โดยไม่ต้องมีเหรียญ}
    \end{choices}
  \end{question}
}

\element{DLT}{
  \begin{question}{QF070}
    Smart Contract มีลักษณะและคุณสมบัติตามข้อใด
    \begin{choices}
        \wrongchoice{สัญญาทางกฎหมายที่พิมพ์ลงในกระดาษชนิดพิเศษ}
        \correctchoice{โปรแกรมคอมพิวเตอร์ที่จะดำเนินการและทำตามเงื่อนไขที่ตกลงไว้โดยอัตโนมัติ}
        \wrongchoice{ซอฟต์แวร์ป้องกันไวรัสที่ฝังตัวในระบบ}
        \wrongchoice{พนักงานกฎหมายที่ทำหน้าที่ตรวจสอบสัญญาผ่านระบบออนไลน์}
    \end{choices}
  \end{question}
}

\element{DLT}{
  \begin{question}{QF071}
    เครือข่ายบล็อกเชนใดที่ได้รับการยอมรับว่าเป็นแพลตฟอร์มแรกที่ริเริ่มนำระบบ Smart Contract มาใช้งานอย่างแพร่หลาย
    \begin{choices}
        \wrongchoice{Bitcoin}
        \correctchoice{Ethereum}
        \wrongchoice{Ripple}
        \wrongchoice{Litecoin}
    \end{choices}
  \end{question}
}

\element{DLT}{
  \begin{question}{QF072}
    Cryptocurrency คือสิ่งใด
    \begin{choices}
        \wrongchoice{เงินกระดาษรูปแบบพิเศษที่พิมพ์ออกมาโดยหน่วยงานเข้ารหัสดิจิทัลของรัฐ}
        \correctchoice{สินทรัพย์ดิจิทัลที่ใช้เทคนิคการเข้ารหัส}
        \wrongchoice{บัตรเครดิตที่ใช้ผูกติดกับบัญชีอีเมลบนอินเทอร์เน็ต}
        \wrongchoice{แอปพลิเคชันสำหรับส่งข้อความแชทที่มีระบบชำระเงิน}
    \end{choices}
  \end{question}
}

\element{DLT}{
  \begin{question}{QF073}
    Node ในเครือข่าย Blockchain หมายถึงอะไร
    \begin{choices}
        \wrongchoice{อุปกรณ์ที่ใช้ขุดบิทคอยน์เท่านั้น}
        \wrongchoice{หน่วยย่อยที่สุดของสกุลเงินดิจิทัลแต่ละประเภท}
        \correctchoice{คอมพิวเตอร์หรืออุปกรณ์ใด ๆ ที่เชื่อมต่อกับเครือข่ายและเก็บรักษาข้อมูลของ Ledger หรือช่วยตรวจสอบธุรกรรม}
        \wrongchoice{กระเป๋าเงินดิจิทัล ส่วนตัวของผู้ใช้งานแบบออฟไลน์}
    \end{choices}
  \end{question}
}

% \element{DLT}{
%   \begin{question}{QF074}
%     Public Blockchain มีลักษณะการทำงานแบบใด
%     \begin{choices}
%         \wrongchoice{จำกัดสิทธิ์เฉพาะพนักงานในบริษัทหรือองค์กรที่กำกับดูแลเท่านั้น}
%         \wrongchoice{ข้อมูลทั้งหมดจะถูกเก็บและควบคุมไว้ที่เซิร์ฟเวอร์หลักขององค์กรผู้สร้างเพียงตัวเดียว}
%         \correctchoice{เปิดกว้างให้ทุกคนสามารถเข้าถึง ทำธุรกรรม และร่วมเป็นโหนดในเครือข่ายได้โดยเสรีไม่ต้องขออนุญาต}
%         \wrongchoice{เป็นบล็อกเชนที่ใช้สำหรับบริการสาธารณะของรัฐบาลโดยมีรัฐบาลเป็นผู้รับรองทุกธุรกรรม}
%     \end{choices}
%   \end{question}
% }

% \element{DLT}{
%   \begin{question}{QF075}
%     Private Blockchain หรือ Permissioned Blockchain มักจะนำไปใช้งานในกรณีใดมากที่สุด
%     \begin{choices}
%         \correctchoice{องค์กรธุรกิจหลายแห่งที่ต้องการแชร์ข้อมูลระหว่างกันในกลุ่มปิดโดยไม่เปิดเผยให้บุคคลทั่วไปภายนอก}
%         \wrongchoice{การสร้างสกุลเงินดิจิทัลระดับโลกที่ต้องการให้ประชาชนทุกคนขุดเก็งกำไรได้}
%         \wrongchoice{เกมออนไลน์แบบเปิดโลกกว้าง (Open World) ที่ต้องการเปิดเสรีไอเทมทุกชิ้นบนตลาดสาธารณะ}
%         \wrongchoice{โซเชียลมีเดียที่เปิดให้ทุกคนอ่านและดูข้อมูลแบบไม่เปิดเผยตัวตน}
%     \end{choices}
%   \end{question}
% }

\element{DLT}{
  \begin{question}{QF076}
    ปัญหาความท้าทายหลักที่สำคัญ อย่างหนึ่งของเทคโนโลยี Blockchain สาธารณะในปัจจุบันคือเรื่องใด
    \begin{choices}
        \wrongchoice{ปัญหาการขาดแคลนภาษาโปรแกรมสำหรับพัฒนาแอปพลิเคชันโดยสิ้นเชิง}
        \wrongchoice{ความเสี่ยงในการคัดลอกไฟล์บัญชีทั้งหมดไปยังแฟลชไดรฟ์โดยแฮกเกอร์ฝีมือดี}
        \correctchoice{ปัญหาความสามารถในการรองรับปริมาณธุรกรรมจำนวนมหาศาลต่อวินาที}
        \wrongchoice{ปัญหาผู้ใช้ทั่วไปลืมรหัสผ่านแล้วทำให้เครือข่ายบล็อกเชนทั้งหมดหยุดชะงัก}
    \end{choices}
  \end{question}
}

\element{DLT}{
  \begin{question}{QF077}
    ข้อใดเป็นตัวอย่างการใช้งาน ของเทคโนโลยี Distributed Ledger Technology (DLT) นอกเหนือจากการเป็นสกุลเงินคริปโต
    \begin{choices}
        \wrongchoice{การพัฒนาโปรแกรมสำหรับตัดต่อวิดีโอที่มีคุณภาพสูง}
        \correctchoice{การใช้ในระบบติดตามห่วงโซ่อุปทาน เพื่อความโปร่งใส}
        \wrongchoice{การช่วยเรนเดอร์กราฟิก 3 มิติเชิงซ้อนในอุตสาหกรรมผลิตภาพยนตร์}
        \wrongchoice{การนำไปใช้เป็นแกนหลักเพื่อเร่งความเร็วระบบปฏิบัติการของสมาร์ทโฟน}
    \end{choices}
  \end{question}
}

\element{DLT}{
  \begin{question}{QF078}
    การโจมตีแบบ 51\% Attack ในระบบเครือข่ายบล็อกเชนหมายถึงสถานการณ์ใด
    \begin{choices}
        \wrongchoice{แฮกเกอร์โจรกรรมเหรียญคิดเป็นมูลค่าลดลงเกินร้อยละ 51 ตลาดคริปโตจึงหยุดชะงัก}
        \wrongchoice{โหนดในระบบเครือข่ายพังพร้อมกันเกินกว่าร้อยละ 51 ส่งผลให้ข้อมูลทั้งหมดหายไปถาวร}
        \correctchoice{ผู้โจมตีสามารถเปลี่ยนแปลงข้อมูลที่บันทึกไว้ในบล็อกเชนได้ มากกว่า 51\% ของเครือข่าย}
        \wrongchoice{ผู้โจมตีเจาะระบบหลอกเอาข้อมูลหรือบัตรประจำตัวของผู้เข้าร่วมจำนวน 51\% ของผู้ใช้งานทั้งหมด}
    \end{choices}
  \end{question}
}

\element{DLT}{
  \begin{question}{QF079}
    Public Key ทำหน้าที่เปรียบเสมือนสิ่งใดในเปรียบเทียบกับระบบการเงินธนาคารแบบดั้งเดิม
    \begin{choices}
        \wrongchoice{รหัสผ่านเข้าใช้แอปพลิเคชันบนมือถือ}
        \correctchoice{หมายเลขบัญชีธนาคาร}
        \wrongchoice{ลายเซ็นส่วนบุคคลด้านหลังเช็คสั่งจ่าย}
        \wrongchoice{ตู้เอทีเอ็มสำหรับทำธุรกรรมสาธารณะ}
    \end{choices}
  \end{question}
}

\element{DLT}{
  \begin{question}{QF080}
    Private Key ของกระเป๋าเงินคริปโตเคอร์เรนซี มีความสำคัญตามข้อใด
    \begin{choices}
        \wrongchoice{ใช้เปิดเผยต่อสาธารณะเพื่อช่วยตรวจสอบความถูกต้องของรายการธุรกรรมที่รับเข้ามา}
        \wrongchoice{ใช้ทำหน้าที่เปรียบเสมือนเครื่องมือค้นหาตำแหน่งของ Block ปัจจุบันในระบบ}
        \correctchoice{เป็นรหัสลับพิเศษที่ใช้ยืนยันการเป็นเจ้าของและใช้ลงนามโอนสินทรัพย์ออกไป}
        \wrongchoice{เป็นเพียงหมายเลขอ้างอิงของการทำธุรกรรมที่สามารถกดขอรับใหม่ผ่านอีเมลได้ หากเกิดลืม}
    \end{choices}
  \end{question}
}

\element{DLT}{
  \begin{question}{QF081}
    ในกรณีที่ผู้ใช้งานทำ Private Key สูญหายและไม่มีคำสำรองไว้ จะเกิดผลลัพธ์ในลักษณะใด
    \begin{choices}
        \wrongchoice{สามารถร้องขอรหัสผ่านใหม่จากฝ่ายสนับสนุน (Support) ของระบบบล็อกเชนได้}
        \wrongchoice{ระบบจะตรวจสอบลายนิ้วมือเพื่อสร้างกุญแจส่วนตัวใหม่ให้โดยอัตโนมัติ}
        \correctchoice{จะไม่สามารถเข้าถึง ควบคุม หรือโอนสินทรัพย์ดิจิทัลเหล่านั้นได้อีกเลย}
        \wrongchoice{ระบบถือว่าบัญชีตาย และสินทรัพย์ทั้งหมดจะถูกโอนกลับไปยังต้นทางของผู้โอนก่อนหน้าอัตโนมัติ}
    \end{choices}
  \end{question}
}

\element{DLT}{
  \begin{question}{QF082}
    กระเป๋าเงินดิจิทัล (Cryptocurrency Wallet) ทำงานโดยมีหลักการหรือหน้าที่หลักคืออะไร
    \begin{choices}
        \wrongchoice{เชื่อมต่อเพื่อทำการขุดเหรียญคริปโต (Mining) ได้โดยตรง}
        \correctchoice{ทำหน้าที่เก็บรักษาและช่วยจัดการการเข้าถึงคู่กุญแจ Public Key และ Private Key อย่างปลอดภัย}
        \wrongchoice{เข้ารหัสข้อมูลไฟล์ภาพถ่ายหรือวิดีโอก่อนอัปโหลดลงบนสื่อสาธารณะ}
        \wrongchoice{ใช้สำหรับตรวจจับไวรัสเรียกค่าไถ่ที่เข้าสู่ระบบคอมพิวเตอร์ผ่านทางอินเทอร์เน็ต}
    \end{choices}
  \end{question}
}

\element{DLT}{
  \begin{question}{QF083}
    ผู้ที่ทำหน้าที่เป็น Miner หรือผู้ขุด ในเครือข่ายแบบ Proof of Work มีบทบาทสำคัญคืออะไร
    \begin{choices}
        \wrongchoice{ทำการลบหรือบีบอัดข้อมูลธุรกรรมเก่าที่ซ้ำซ้อนเพื่อเพิ่มเนื้อที่จัดเก็บในระบบ}
        \correctchoice{รวบรวมและตรวจสอบธุรกรรม แล้วแข่งกันแก้โจทย์คณิตศาสตร์ที่ซับซ้อนเพื่อสร้างบล็อกใหม่}
        \wrongchoice{เขียนโปรแกรมสำหรับการสร้างเหรียญชนิดใหม่เตรียมเข้าสู่ตลาดแลกเปลี่ยนคริปโต}
        \wrongchoice{ร่างกฎหมายและนโยบายกำกับดูแลเสถียรภาพมูลค่าเหรียญเพื่อให้มีราคาคงที่}
    \end{choices}
  \end{question}
}
